\documentclass[a4paper,11pt]{article}
\usepackage[a4paper, margin=2cm]{geometry}
\usepackage[T1]{fontenc}

\usepackage[backend=biber]{biblatex}
\addbibresource{bibliography.bib}

\usepackage{
    amsmath,amssymb,fontawesome,graphicx,mathpartir,
    mathtools,minted,simplebnf,stmaryrd,xcolor
}

\newcommand{\meet}[2]
{\rangle\joinrel\joinrel\overset{#1}{
    \underset{#2}{\sim}}\joinrel\joinrel\langle}

\usepackage{amsthm}
\newtheorem{theorem}{Theorem}
\theoremstyle{definition}
\newtheorem{definition}{Definition}

\usepackage{newunicodechar}
\newunicodechar{∀}{$\forall$}
\newunicodechar{∷}{$\dblcolon$}

\newcommand{\SSTLC}{\llbracket\mathsf{S}\rrbracket\mathsf{TLC}}
\newcommand{\twoLTT}{\llbracket 2\rrbracket\mathsf{LTT}}

\usepackage{authblk,orcidlink}

\title{Towards Type Systems That Interpret Themselves}

\author[ ]{Pavel Sokolov\raisebox{-0.45ex}{\scriptsize\orcidlink{0009-0001-6922-2870}}\kern-0.4em\textsuperscript{1} \guilsinglleft\protect\href{mailto:ppsokolov@hse.ru}{\texttt{ppsokolov@hse.ru}}\guilsinglright
\authorcr
Alexander Kuklev\raisebox{-0.45ex}{\scriptsize\orcidlink{0009-0004-1019-4205}}\kern-0.4em\textsuperscript{2} \guilsinglleft\protect\href{mailto:a@kuklev.com}{\texttt{a@kuklev.com}}\guilsinglright}

\affil[1]{\small HSE University, Laboratory of
Theoretical Computer Science}

\affil[2]{Radboud Univ. Nijmegen, Department of
Software Science}

\begin{document}

\maketitle

\begin{abstract}
    We search for type systems capable of ``eating
    itself'' in the sense that their own metatheoretical
    properties could be stated and proven inside it. In
    this particular instance, we study consequences of
    introducing an interpretation operator $m\llbracket
    t\rrbracket$ which allows instantiation of a closed
    term $t$ as a construction inside a model $m$
    definable internally. Our main result is a design of
    a two-level type theory with interpretation,
    $\llbracket 2\rrbracket\mathrm{LTT}$, where
    intrepretation can be used to provide a limited form
    of universe polymorphism, prove external
    parametricity and do some synthetic mathematics.
\end{abstract}

\section{Introduction}

An idea that a ``type theory should eat itself'' is a
well-established concept in study of dependent type
theory. Under this guise, if a type theory is presented
as a ``full-scale system for constructive mathematics``
[CITE], then it should be capable of proving its own
metatheoretical properties. For example, in a paper
bearing the same name [CITE], the author presents a
normalizer for type theory written in Agda [CITE].

Of course, G\"odel's incompleteness theorems [CITE]
still apply in a constructive context, so na\"ively
proving self-consistency of a type system supporting
primitive recursion is out of question. What is usually
proven instead is that a particular implementation of a
type theory inside itself is sound assuming the
consistency of the theory, so no paradox arises [CITE]
[CITE]; this is still useful because it leads to the
phenomena of \emph{certified proof assistants}.

Another possible approach is to work in a two-level type
theory [CITE]. As the name implies, such a theory
consists of two levels: the lower \emph{object} level is
the main type theory, and a higher \emph{meta} level is
used to internalize statements about object level which
are external to it. A classical example given in the
cited paper is a construction of a \emph{semisimplicial
type of level $n$} for each external $n\in\mathbb{N}$;
object level here is homotopy type theory (HoTT) with
univalence axiom (UA) [CITE] and meta level is
Martin-L\"of type theory (MLTT) with uniqueness of
identity proofs (UIP) [CITE]; in usual type theory,
these two axioms are mutually exclusive, but the trick
is that an object-level equality becomes a different
type when viewed from meta level, so UA and UIP apply to
different notions of ``equality'', keeping the system
consistent (and sound with existing models of HoTT
[CITE]).

We aim to generalize this construction in the following
way. It can be said that an object level has a fixed
\emph{interpretation} in a meta-level; what if the
system allowed to vary the interpretation? In other
words, what if the system included a generic
interpretation operator denoted $m\llbracket t
\rrbracket$, where $t$ is a term which is being
interpreted, and $m$ is a model definable inside the
system? (We will call such systems \emph{reflective}.)

Turns out, we would get quite a few useful
metatheoretical results about object-level theory for
free; in particular, Wadler's parametricity theorem
[CITE] (for its formulation in dependently typed
context, see [CITE] and [CITE]) about object-level
theory would be provable on the meta-level by means of
specifying a relational model; a fraction of displayed
type theory [CITE] can be obtained via similar
translation of a type into a unary predicate. Another
fun consequence is that object-level theory would get
universe polymorphism via displacements as described by
Conor Titania McBride [CITE] and Favonia [CITE]; we
suspect that there might be other useful applications of
internal interpretation.

\subsection{Contributions}

So the goal is to construct type systems with
        such interpretation present as ; we will call type system with such
interpretation capabilities \emph{reflective}.
\begin{itemize}
    \item So our main goal is to construct an
        implementable reflective type system. The main
        difficulty in constructing such a system arises
        from a certain no-go result discussed in
        Section~\ref{no-go-section}, which forces us to
        limit the power of reinterpretation in order to
        keep type checking decidable.
    \item For ease of presentation, we first offer a
        reflective type system based on the
        Hindley-Milner type system, where object level
        is simply-typed lambda calculus with product
        types (STLC), called $\SSTLC$
        (Section~\ref{HM-R}). We prove that it is
        conservative with respect to STLC (hence
        strongly normalizing and sound), and that global
        type inference is decidable (section
        \ref{HM-R-properties}).
    \item Building upon the insights gathered from
        developing $\SSTLC$, we are finally equipped to
        dive into the reflective version of two-level
        type theory, where object level is vanilla MLTT,
        dubbed $\twoLTT$ (Section~\ref{MLTT-R}). We
        prove soundness and conservativity of it as well
        (Section~\ref{MLTT-R-properties}).
    \item As intended, by virtue of reinterpretation,
        $\twoLTT$ incorporates both a restricted form of
        universe polymorphism and external
        parametricity. And, of course, you can use it as
        a system for synthetic mathematics with internal
        reuse of synthetic results
        (Section~\ref{MLTT-R-applications}).
\end{itemize}

\section{No-go theorem of internal interpretation}
\label{no-go-section}

So we want to build a reflective type system. That is,
a system where a user is able to both define a model of
a system internally and apply it to well-typed terms to
get their internal interpretations. Naturally, we would
like user-defined models to respect definitional
equalities of a system; in a dependently typed system,
this could be naturally achieved by using an identity
type to spell out corresponding equalities inside a
model. Then, it would be nice if, before interpreting
a term, system would first bring it to normal form to
get rid of the artifacts of the construction in the
result of the interpretation. Finally, in a dependently
typed system, it is natural to expect reductions of open
terms, so we would like our interpretation to compute on
open terms.

The theorem proven below asserts that we cannot have all
three in a single system.

First, a couple of definitions. We're going to work
within a classical type-theoretical framework, which
introduces a sort of contexts $\Gamma$, terms $t,t'$ and
types $T$ linked by the typing relation
$\Gamma\vdash t:T$; computation is specified by a
one-step reduction relation $t\leadsto t'$ on terms.

\begin{definition}
    [Na\"ive computational interpretation]
    We call a type operator $F[\_]$ a na\"ive
    computational interpretation of types if it
    respects computation on types (that is, if there is
    a reduction relation on types $(\leadsto)$ and
    $T\leadsto T'$, then $F[T]\leadsto^* F[T']$), and
    computes as a primitive recursive function on normal
    forms of types. Analogously, term with a hole
    $\llbracket\_\rrbracket$ is a na\"ive
    computational interpretation of terms if it
    respects computation on terms, computes as a
    primitive recursive function on normal forms of
    terms and, if $\Gamma\vdash t:T$, then
    $\Gamma\vdash\llbracket t\rrbracket:F[T]$.
\end{definition}

\begin{theorem}[No-go theorem of naive interpretation]
\label{no-go}
    In a confluent type system with a type of functions
    and na\"{\i}ve computational interpretation, if
    reduction commutes with substitution, for arbitrary
    well-typed $s$ and $t$ there exists $u$ such that:
    \[
        \mathtt{app}\llbracket\lambda x.t\rrbracket
        \llbracket s \rrbracket \leadsto^* u\;
        \text{\reflectbox{$\leadsto^*$}}\;
        \llbracket t[s/x] \rrbracket
    \]
    where $\mathtt{app}$ is how $\llbracket\_\rrbracket$
    acts on applications.
\end{theorem}

Therefore, $\mathtt{app}\llbracket\lambda x.t\rrbracket
\llbracket s \rrbracket$ and $\llbracket t[s/x]
\rrbracket$ are definitionally equivalent, which forces
beta-reductions in a model to hold definitionally rather
than propositionally, directly contradicting our desired
outcome. Analogous theorems can be proven for any other
reduction present in the system; in other words,
\textbf{propositional equality is not strong enough}.
The proof is by direct example:

\begin{proof}
    We are given well-typed $\Gamma\vdash s:S$ and
    $\Gamma,x:S\vdash t:T$. Consider the following term:
    \[
        \Gamma,x:S\to T,y:S\vdash
            \llbracket x\,y\rrbracket:F[T]
    \]
    Term under interpretation is in normal form, hence
    the whole term computes to $\mathtt{app}\llbracket x
    \rrbracket\llbracket y\rrbracket$.
    Substituting $\lambda x.t$ for $x$ and $s$ for $y$,
    we get $\mathtt{app}\llbracket\lambda x.t\rrbracket
    \llbracket s \rrbracket$.

    On the other hand, if we perform substitution first,
    we would get $\Gamma\vdash\llbracket(\lambda x.t)s
    \rrbracket:F[T]$. Since na\"{\i}ve interpretation
    respects reduction, this reduces to
    $\llbracket t[s/x] \rrbracket$.

    Since reduction commutes with substitution and
    system is confluent, there should be such $u$ that
    both terms reduce to.
\end{proof}

Fundamentally, the proof above exploits the openness of
the original term; if we restrict interpretation to
compute on closed terms only, we can arrive at a sound
system without requiring definitional equalities in
models, as exemplified by our non-dependently typed
system $\SSTLC$ described in Section~\ref{SSTLC}. And
this, actually, is a well-known setup: under the same
conditions (no dependent types + computation on closed
terms), self-interpreter for System~F$\omega$ was
developed \cite{brown2016breaking}.

For a dependently typed system, however, our no-go
theorem leaves us with three options:

\begin{enumerate}
    \item Restrict models to those where equalities hold
        definitionally rather than propositionally.
    \item Do not proceed with interpretation unless
        the term is closed. To allow reasoning on open
        terms, equations describing the behavior of
        interpretation are exposed as propositional
        axioms.
    \item Do not reduce open terms under
        interpretation at all, performing the
        interpretation as-is. To simplify the result
        of interpretation, one can rewrite it with
        equalities present in the model.
\end{enumerate}

If you view existing work on displacement
\cite{McBride:Crude2011, McBride:Crude2012}, internal
parametricity \cite{Altenkirch:Parametricity2024} and
displayed type theory \cite{Kolomatskaia_Shulman_2025}
through the lens of interpretation inside a fixed
model (or family of models), you'll see that
corresponding models indeed validate equalities
definitionally; thus it seems that the first approach
needs to be taken. It is also the most satisfying from
the computational point of view. However, if you want a
sound system where models can be user-defined, you need
to constrain models through typing discipline. By
following techniques currently known to us
\cite{bezem2021notegeneralizedalgebraictheories},
constraining models in such a way would require
defining a type of models as a type of algebras of an
inductive-inductive-recursive type family (or even
inductive-inductive-recursive simplicial types, for
univalent type theories) which is well beyond the scope
of this paper.

Out of the two remaining options, the last one is better
suited for practical dependently-typed systems since it
properly computes on open terms. In addition, it doesn't
restrict the class of possible models while at the same
time providing the benefits of the first approach: if
a model is such that the equations hold definitionally,
no rewriting on the result of interpretation is needed
to reduce it. The only downside is that, possibly,
reductions \textit{after} interpretation are more costly
than the ones we could perform \textit{before} it.

\section{
    Almost simply typed almost reflective language:
    \texorpdfstring{$\SSTLC$}{[[S]]TLC}
}\label{SSTLC}

We start with a simpler language to iron out the details
and actually figure out how to properly reinterpret
terms internally before moving on to a dependently typed
system. The simplest language for which we know a lot of
interesting models is the simply typed lambda calculus
with products, which we refer to as ``STLC'' throughout
this paper; in this section, we'll present a language
capable of interpreting STLC terms as constructions in
STLC models that can be defined in our language.

So, what are the models of STLC? In addition to the
usual set-theoretic semantics where $S\to T$ is
interpreted as a function from $S$ to $T$ and
$S\times T$ is a cartesian product, there are
proof-theoretic semantics that relate terms of type $T$
to intuitionistic proofs of proposition
$\llbracket T\rrbracket$, where $(\to)$ is now
understood to mean implication, and $(\times)$ is
conjunction. This, of course, is a classic Curry-Howard
correspondence \cite{ho80}; both set-theoretic and
proof-theoretic semantics are generalized by a complete
description of STLC models as given by Joachim Lambek
\cite{Lambek1980}. According to this
\textit{Curry-Howard-Lambek} correspondence, models
of STLC are precisely the cartesian closed categories
(or CCCs for short). Definitions vary; we're going to
use the definition \ref{ccc}.

\begin{definition}[Cartesian closed category]\label{ccc}
A cartesian closed category is the following datum:
\begin{enumerate}
\item A class\footnote{Objects do not necessarily form a
    set: for example, one might want to have a category
    of sets where morphisms are functions between them.
    To deal with sizing issues, pick appropriate
    set-theoretical foundations you are most comfortable
    with (NBG, MK, ZFC/S or something else); this detail
    is not important in our development} of
    \textit{objects} $\mathrm{Ob}$ and a family of
    \textit{morphisms} $\mathrm{Hom}(S, T)$ indexed by
    $S,T\in\mathrm{Ob}$, where $f\in\mathrm{Hom}(S,T)$
    is denoted by $f:S\to T$. In addition, $S$ is called
    the \textit{source} of $f$, and $T$ is the
    \textit{target} of $f$.
\item For each $A\in\mathrm{Ob}$ there's an
    \textit{identity morphism} $\mathrm{id}_A: A\to A$;
    for each $f:A\to B$ and $g:B\to C$, there's their
    \textit{composition} $g\circ f:A\to C$. If a
    source-target of an identity morphism is clear from
    the context, the subscript is elided. Also, the
    following properties hold:
    \begin{align*}
        \mathrm{id}\circ f = f = f\circ\mathrm{id} & &
        (f \circ g) \circ h = f \circ (g \circ h)
    \end{align*}
    This applies whenever all compositions are well-defined.
\item There's a \textit{terminal object}
    $\top\in\mathrm{Ob}$ such that for each
    $A\in\mathrm{Ob}$ there's a \textbf{unique} morphism
    $!_A:A\to\top$.
\item For each $A,B\in\mathrm{Ob}$ there's their
    \textit{product} $A\times B$ together with
    \textit{projections} $\pi_1:A\times B\to A$ and
    $\pi_2:A\times B\to B$; for each object
    $C\in\mathrm{Ob}$ and a pair of morphisms $f:C\to A$
    and $g:C\to B$, there's a \textbf{unique}
    \textit{pairing morphism} $[f;g]:C\to A\times B$
    such that the following properties hold:
    \begin{align*}
        \pi_1\circ[f;g]=f & & \pi_2\circ[f;g]=g
    \end{align*}
\item For each $A,B\in\mathrm{Ob}$ there's their
    \textit{exponential object} $B^A$ together with
    \textit{evaluation morphism}
    $\mathrm{ev}:B^A\times A\to B$; for each object
    $C\in\mathrm{Ob}$ and morphism $f:C\times A\to B$,
    there's a \textbf{unique} \textit{abstraction}
    $\lambda\,f:C\to B^A$ such that the following
    property holds for arbitrary $g:C\to A$:
    \[
        \mathrm{ev}\circ[\lambda\,f; g]
        = f\circ[\mathrm{id}_C; g]
    \]
\end{enumerate}
\end{definition}

Then, a well-typed term $\Gamma\vdash t:T$ is
interpreted as a morphism $\llbracket t\rrbracket:
\llbracket\Gamma\rrbracket\to\llbracket T\rrbracket$,
where $\llbracket\Gamma\rrbracket$ is a left-nested
product of interpretations of types from $\Gamma$:
\[
    \llbracket x:A,y:B,z:C\rrbracket=
    ((\top\times\llbracket A\rrbracket)
    \times\llbracket B\rrbracket)
    \times\llbracket C\rrbracket.
\]
Unit type is now a terminal object; a product of types is
now a product of objects; an arrow type is an exponential
object. Equations in (2) are needed to validate the
structural properties of substitution; equations in (4)
correspond to reductions of pairs; equation in (5)
validates $\beta$-reduction. \textbf{Unique}ness
properties correspond to $\eta$-equivalence.

For example, in our language, $m\,\llbracket t
\rrbracket$ is the operator interpreting the term $t$ in
an internal model $m$. On a sample term
\[
    \cdot\vdash\lambda f x y. (f x, y):
        (A\to B)\to A\to C\to B\times C
\]
it computes like this:
\begin{multline*}
    m\,\llbracket\lambda f x y. (f x, y))\rrbracket=\\=
    \mathtt{lam}\,(\mathtt{lam}\,(\mathtt{lam}\,(
        \mathtt{pair}\,
            (\mathtt{eval}\,\circ\,\mathtt{pair}\,
                (\mathtt{proj}_2\circ\mathtt{proj}_1
                    \circ\mathtt{proj}_1
                )\,
                (\mathtt{proj}_2\circ\mathtt{proj}_1)
            )\,
            \mathtt{proj}_2
    )))
\end{multline*}

Where $\mathtt{lam}$,\,$\mathtt{pair}$ and
$\mathtt{eval}$ are the abstraction morphism, the
pairing morphism, and the evaluation morphism taken from
$m$, defined in the same way as in the definition
\ref{ccc}, translated into types of our system.

Now, moving on to the actual design of our language. As
discussed in Section~\ref{no-go-section}, for the
non-dependent case we are now dealing with, we will
interpret only normal forms of closed terms. Among other
things, this means that we actually don't need equations
and \textbf{unique}nesses from Definition~\ref{ccc} to
hold in our internal models! This approach doesn't break
any of the metatheoretical properties considered in
Section~\ref{SSTLC-properties}. Additionally, for the
sake of example, we're only going to consider models
where objects are types. Then, the internalization of
definition \ref{ccc} would require 1) type operators to
define types of morphisms, products, and exponential
objects; 2) polymorphic functions to operationalize the
$\forall\exists$ introductions of constructions in a
model. Actually, we're going to get rid of type
operators as well, since we can do with just
types-in-a-context which are then substituted into
the definition of a model for checking term-level
operations. This way, our language is a small extension
of the classic Hindley-Milner type system; grammar rules
are collected in a Figure~\ref{sstlc-grammar}. You'll
notice that, despite our initial description, grammar
doesn't define a reinterpretation operator
$m\llbracket t\rrbracket$ for arbitrary term $t$;
instead, only variables (denoted by a metavariable $x$)
get to be reinterpreted. In this system,
reinterpretation is nothing but a total function on
closed values which is called whenever a value is
substituted into such an $x$ (so, in a way, this is
similar to hereditary substitution
\cite{watkins2003concurrent,keller2010hereditary}).

\begin{definition}[Values of $\SSTLC$]
Values are denoted by metavariables $u,v$ in the BNF
grammar below. Metavariable $n$ denotes neutral terms:
\begin{align*}
    u,v&\Coloneqq n\;|\;()\;|\;(u,v)\;|\;\lambda x.v\\
    n&\Coloneqq x\;|\;\pi_i\,n\;|\;n\,v
\end{align*}
\end{definition}

\begin{figure}
\centering
\begin{bnf}
    $S,T,C,P,E,TT$ : \textsf{MonoType} ::=
    | $\alpha$ : type variable
    | $\top$ : unit type
    | $S\times T$ : product type
    | $S\to T$ : arrow type
    ;;
    $\Delta$ : \textsf{MonoCtx} ::=
    | $\cdot$ : empty context
    | $\Delta,x\colon T$ : context extension
    ;;
    $\tau$ : \textsf{PolyType} ::=
    | $T$ : monomorphic type
    | $\forall\alpha.\tau$ : forall type
    ;;
    $\Gamma$ : \textsf{PolyCtx} ::=
    | $\cdot$ : empty context
    | $\Gamma,x\colon\tau$ : context extension
    ;;
    $\Xi$ : \textsf{Ctx} ::=
    | $\Gamma\,\|\,\Delta$ : joined context
    ;;
    $s,t$ : \textsf{Term} ::=
    | $x$ : variable
    | $()$ : singleton
    | $(s, t)$ : pair
    | $\pi_1\,t$ : first projection
    | $\pi_2\,t$ : second projection
    | $\lambda x. t$ : abstraction
    | $t\,s$ : application
    | \texttt{let}\,$x=s$\,\texttt{in}\,$t$
        : polymorphic binding
    | $m\llbracket x\rrbracket$ : reinterpretation
    ;;
    $m$ : \textsf{Model} ::=
    | $\langle T, t, \ldots \rangle$
        : model def. (see typing rules)
\end{bnf}
\caption{$\SSTLC$ terms and values}
\label{sstlc-grammar}
\end{figure}

Thus, to define interpretation, it is enough to consider
STLC terms:

\begin{definition}[Interpretation of terms in $\SSTLC$]
    Given a model $m$, a list of variables $\gamma$ and
    a term $t$, interpretation of a term $m\llbracket
    \gamma\vdash t\rrbracket$ is a recursive function
    defined by the following equations:
    \begin{align}{}
        m\llbracket\gamma\vdash()\rrbracket&=
        \mathtt{cnst}
        \\
        m\llbracket\gamma\vdash(s,t)\rrbracket&=
        \mathtt{pair}
        \,m\llbracket\gamma\vdash s\rrbracket
        \,m\llbracket\gamma\vdash t\rrbracket
        \\
        m\llbracket\gamma\vdash\pi_i\,t\rrbracket&=
        \mathtt{comp}\,\mathtt{proj}_i\,
        m\llbracket\gamma\vdash t\rrbracket
        \\
        m\llbracket\gamma\vdash x\rrbracket&=\gamma[x]
        \\
        m\llbracket\gamma\vdash\lambda x.t\rrbracket&=
        \mathtt{lam}\,
        m\llbracket(\gamma,x)\vdash t\rrbracket
        \\
        m\llbracket\gamma\vdash t\,s\rrbracket&=
        \mathtt{comp}\,\mathtt{eval}\,\left(
        \mathtt{pair}
        \,m\llbracket\gamma\vdash t\rrbracket
        \,m\llbracket\gamma\vdash s\rrbracket
        \right)
    \end{align}
    where model operations \texttt{cnst},\,\texttt{pair}
    and others are taken from $m$ and $\gamma[x]$ is
    also defined by equations:
    \begin{align}
        (\gamma,x)[x] &= \mathtt{proj}_2 \\
        (\gamma,y)[x] &= \mathtt{comp}\,\gamma[x]\,
        \mathtt{proj}_1
    \end{align}
\end{definition}

Here is our way to enforce a principle that \textit{only
closed terms get interpreted}: once a term $s$ bound by
a polymorphic let-binding
$\mathtt{let}\,x=s\,\mathtt{in}\,t$ becomes closed and
is brought into normal form, only then is it
substituted into $t$, and whenever it is substituted
into $m\llbracket x\rrbracket$, it immediately gets
reinterpreted. This operation is, as noted before,
similar to hereditary substitution; it is denoted by
$t[s\sslash x]$ in our computation rules given in
Figure~\ref{sstlc-comp}. So it acts like a regular
substitution in all cases except for variables under
reinterpretation:
\[
    m\llbracket x\rrbracket[s\sslash x] =
    m\llbracket()\vdash s\rrbracket
\]
One can easily notice that computation would get stuck
if let-bound terms were to contain variables coming
from unapplied $\lambda$-abstractions. Our typing rules
given in Figure~\ref{sstlc-typing} protect us from that,
allowing to refer only to let-bound terms inside other
let-bound terms. Again, aside from that and the typing
rule for universes, there's nothing different from the
usual declarative typing rules for a Hindley-Milner type
system \cite{Clement:ML1986,damas1982principal,
milner1978theory}.

\begin{figure}
\centering
\begin{mathpar}
    \inferrule{s\leadsto s'}{(s,t)\leadsto(s',t)}
    \and
    \inferrule{t\leadsto t'}{(s,t)\leadsto(s,t')}
    \and
    \inferrule{t\leadsto t'}{\pi_i t\leadsto\pi_i t'}
    \and
    \inferrule{}{\pi_i(t_1, t_2)\leadsto t_i}
    \\
    \inferrule{t\leadsto t'}
        {\lambda x.t\leadsto\lambda x.t'}
    \and
    \inferrule{s\leadsto s'}{s\,t\leadsto s'\,t}
    \and
    \inferrule{t\leadsto t'}{s\,t\leadsto s\,t'}
    \and
    \inferrule{}{(\lambda x.t)s\leadsto t[s/x]}
    \\
    \inferrule{s\leadsto s'}{
        \mathtt{let}\,x=s\,\mathtt{in}\,t\leadsto
        \mathtt{let}\,x=s'\,\mathtt{in}\,t
    }
    \and
    \inferrule{t\leadsto t'}{
        \mathtt{let}\,x=s\,\mathtt{in}\,t\leadsto
        \mathtt{let}\,x=s\,\mathtt{in}\,t'
    }
    \and
    \inferrule{\mathrm{fv}(v)=\varnothing}{
        \mathtt{let}\,x=v\,\mathtt{in}\,t
        \leadsto t[v\sslash x]
    }
\end{mathpar}
\caption{$\SSTLC$ computation rules}
\label{sstlc-comp}
\end{figure}

To check models, we abuse our type system a bit 
to avoid adding type operators and complicating it:
types of morphisms $C$, product objects $P$ and
exponentials $E$ are specified as open monomorphic types
which can use two type variables, $0$ and $1$. This fact
is reflected in the syntax of a model using notation
$\Lambda 01.C$, $\Lambda 01.P$ and $\Lambda 01.E$
respectively. $0$ and $1$ are then filled in by
substitution when we need to type model operations
(notation $T[S_1, S_2]$ is a shorthand for
$T[S_1/0, S_2/1]$). The terminal object is
just a type named $TT$.

In order to check reinterpretation, we first need to
define how the model acts on monomorphic types. As there
is no computation on them, reinterpretation of types is
defined once again as a recursive function on syntax.

\begin{definition}[Interpretation of types in $\SSTLC$]
    Let $m$ be a correct model and $\sigma$ be a
    substitution of monomorphic types. Then, the
    interpretation of a type in a model is the following
    function:
    \begin{align}{}
        m\llbracket\sigma\vdash\alpha\rrbracket&=
            \sigma\alpha \\
        m\llbracket\sigma\vdash\top\rrbracket&=TT \\
        m\llbracket\sigma\vdash(S\times T)\rrbracket&=
            P[m\llbracket\sigma\vdash S\rrbracket,
            m\llbracket\sigma\vdash T\rrbracket] \\
        m\llbracket\sigma\vdash S\to T\rrbracket&=
            E[m\llbracket\sigma\vdash S\rrbracket,
            m\llbracket\sigma\vdash T\rrbracket].
    \end{align}
\end{definition}

\begin{figure}
\centering
\begin{mathpar}
    \inferrule{}{\Xi\vdash ():\top}
    \and
    \inferrule
    {\Xi\vdash s:S\\\Xi\vdash t:T}
    {\Xi\vdash (s,t):S\times T}
    \and
    \inferrule
    {\Xi\vdash t:S\times T}
    {\Xi\vdash \pi_1\,t:S}
    \and
    \inferrule
    {\Xi\vdash t:S\times T}
    {\Xi\vdash \pi_2\,t:T}\\
    \and
    \inferrule{x:T\in\Delta}
    {\Gamma\,\|\,\Delta\vdash x:T}
    \and
    \inferrule{\Gamma\,\|\,\Delta,x:S\vdash t:T}
        {\Gamma\,\|\,\Delta\vdash\lambda x.t:S\to T}
    \and
    \inferrule
    {\Xi\vdash t:S\to T\\\Xi\vdash s:S}
    {\Xi\vdash t\,s:T}
    \and
    \inferrule{
        \Gamma\vdash s:S\\
        \mathrm{fv}(S)=\vec{\alpha}\\
        \Gamma,x:\forall\vec{\alpha}.S\,\|\,\Delta
        \vdash t:T
    }{\Gamma\,\|\,\Delta\vdash
        \mathtt{let}\,x=s\,\mathtt{in}\,t:T}
    \and
    \inferrule{
        \Gamma\,\|\,\Delta\vdash m\,\mathrm{model}\\
        x:\forall\vec{\alpha}.T\in\Gamma\\
        \sigma=\left[\vec{S}/\vec{\alpha}\right]
    }{\Gamma\,\|\,\Delta\vdash m\llbracket x\rrbracket
        :C_m[TT_m, m\llbracket\sigma\vdash T\rrbracket]}
    \and
    \inferrule{
        \mathrm{fv}(C)\cup\mathrm{fv}(P)\cup
            \mathrm{fv}(E)\subseteq\{0,1\}\\
        \Xi\vdash\mathtt{id}:C[\alpha,\alpha]\\
        \Xi\vdash\mathtt{comp}:
            C[\beta,\gamma]\to C[\alpha,\beta]\to
            C[\alpha,\gamma]\\
        \mathrm{fv}(TT)=\varnothing\\
        \Xi\vdash\mathtt{cnst}:C[\alpha,TT]\\
        \Xi\vdash\mathtt{pair}:
            C[\gamma,\alpha]\to C[\gamma,\beta]\to
            C[\gamma,P[\alpha,\beta]]\\
        \Xi\vdash\mathtt{proj}_1:
            C[P[\alpha,\beta],\alpha]\\
        \Xi\vdash\mathtt{proj}_2:
            C[P[\alpha,\beta],\beta]\\
        \Xi\vdash\mathtt{lam}:
            C[P[\gamma,\beta],\alpha]\to
            C[\gamma,E[\alpha,\beta]]\\
        \Xi\vdash\mathtt{eval}:
            C[P[E[\alpha,\beta],\beta],\alpha]
    }{\Xi\vdash\langle
        \Lambda 01.C,\mathtt{id},\mathtt{comp}, TT,
        \mathtt{cnst}, \Lambda 01.P,\mathtt{pair},
        \mathtt{proj}_1,\mathtt{proj}_2,
        \Lambda 01.E,\mathtt{lam},\mathtt{eval}
    \rangle\,\mathrm{model}}
\end{mathpar}
\caption{$\SSTLC$ declarative typing rules}
\label{sstlc-typing}
\end{figure}

\subsection{Metatheoretic properties of
\texorpdfstring{$\SSTLC$}{[[S]]TLC}}
\label{SSTLC-properties}

Here we only study some basic properties of a system to
make sure that everything works as expected:
computations are confluent; types are preserved by all
computations; all the types can be inferred if you
wish to do so; everything terminates, and, modulo
reduction of \texttt{let}-bindings, this is the same old
STLC we started with. Proofs mostly follow
standard techniques since the rules for reinterpretation
are structured to directly plug into proofs of
properties of the STLC and the Hindley-Milner type
system.

\begin{theorem}[Church-Rosser theorem for $\SSTLC$]
    Reflexive transitive closure of $(\leadsto)$ has
    the Church-Rosser property.
\end{theorem}
\begin{proof}
    Classic proof via parallel reduction relation
    \cite{takahashi1995parallel,barendregt1984lambda}.
    Parallel reduction is defined as follows:
    \begin{mathpar}
        \inferrule{}{t\rightrightarrows t}
        \and
        \inferrule{
            s\rightrightarrows s'\\
            t\rightrightarrows t'
        }{(s,t)\rightrightarrows(s',t')}
        \and
        \inferrule{t\rightrightarrows t'}
            {\pi_i\,t\rightrightarrows\pi_i\,t'}
        \and
        \inferrule{t_1\rightrightarrows t_1'}
            {\pi_1(t_1,t_2)\rightrightarrows t_1'}
        \and
        \inferrule{t_2\rightrightarrows t_2'}
            {\pi_2(t_1,t_2)\rightrightarrows t_2'}
        \\
        \inferrule{t\rightrightarrows t'}
            {\lambda x.t\rightrightarrows\lambda x.t'}
        \and
        \inferrule{
            s\rightrightarrows s'\\
            t\rightrightarrows t'
        }
        {s\,t\rightrightarrows s'\,t'}
        \and
        \inferrule{
            s\rightrightarrows s'\\
            t\rightrightarrows t'
        }
        {(\lambda x.s)\,t\rightrightarrows
            s'[t'/x]}
        \\
        \inferrule{
            s\rightrightarrows s'\\
            t\rightrightarrows t'
        }{
            \mathrm{let}\,x=s\,\mathrm{in}\,t
                \rightrightarrows
            \mathrm{let}\,x=s'\,\mathrm{in}\,t'
        }
        \and
        \inferrule{t\rightrightarrows t'}{
            \mathrm{let}\,x=v\,\mathrm{in}\,t
                \rightrightarrows
            t'[v\sslash x]
        }
    \end{mathpar}
\end{proof}

\begin{theorem}[Type preservation for $\SSTLC$]
    The following is true:
    \begin{enumerate}
        \item If $\Xi\vdash m\,\mathrm{model}$,
            $\Delta\vdash t:T$ and $\delta$ is a list of
            variables from $\Delta$, then \\ $\Xi\vdash
            m\llbracket\delta\vdash t\rrbracket:
            C_m[m\llbracket\Delta\rrbracket,
            m\llbracket T\rrbracket]$.
        \item If $\Gamma\,\|\,\Delta\vdash s:S$ and
            $\Gamma\,\|\,\Delta,x:S,\Delta'\vdash t:T$,
            then $\Gamma\,\|\,\Delta,\Delta'\vdash
            t[s/x]:T$.
        \item If $\Gamma\vdash v:\tau$ and
            $\Gamma,x:\tau,\Gamma'\,\|\,\Delta\vdash
            t:T$, then
            $\Gamma,\Gamma'\,\|\,\Delta\vdash
            t[v\sslash x]:T$.
        \item If $\Gamma\vdash t:T$ and $t\leadsto t'$,
            then $\Gamma\vdash t':T$.
    \end{enumerate}
\end{theorem}
\begin{proof}
    By structural induction on terms and $(\leadsto)$.
\end{proof}

\begin{theorem}[Strong normalization for $\SSTLC$]
    Closed terms typable in $\SSTLC$ are strongly
    normalizing with respect to $(\leadsto)$ and reduce
    to a value.
\end{theorem}
\begin{proof}
    Classic proof via saturated sets
    \cite{tait1967intensional,girard1989proofs,
    luo1994computation}. Let $\mathsf{SN}$ be a set of
    terms which are strongly normalizing with respect to
    $(\leadsto)$ and normal forms of which are values.
    Let $\mathsf{Atomic}$ be the least set of terms
    satisfying the following conditions: variables $x$
    are $\mathsf{Atomic}$; if $t$ is $\mathsf{Atomic}$,
    then so are $\pi_i\,t$ and $t\,s$. Then, let
    $t\to_\mathrm{WHR}t'$ be the following relation on
    terms:
    \begin{mathpar}
        \inferrule{t\in\mathsf{SN}}
        {(\lambda x.s)t\to_\mathrm{WHR}s[t/x]}
        \and\inferrule{t_2\in\mathsf{SN}}
        {\pi_1(t_1,t_2)\to_\mathrm{WHR}t_1}
        \and\inferrule{t_1\in\mathsf{SN}}
        {\pi_2(t_1,t_2)\to_\mathrm{WHR}t_2}
        \and\inferrule{t\to_\mathrm{WHR}t'}
        {t\,s\to_\mathrm{WHR}t'\,s}
        \and\inferrule{t\to_\mathrm{WHR}t'}
        {\pi_i\,t\to_\mathrm{WHR}\pi_i\,t'}
        \and\inferrule{ }
        {\mathtt{let}\,x=v\,\mathtt{in}\,t
        \to_\mathrm{WHR}t[v\sslash x]}
        (\mathrm{fv}(v)=\varnothing)
    \end{mathpar}

    Finally, let $\mathsf{SAT}$ be the following family
    of sets of terms:
    \begin{itemize}
        \item if $X\in\mathsf{SAT}$, then
            $X\subseteq\mathsf{SN}$;
        \item if $X\in\mathsf{SAT}$, then
            $\mathsf{Atomic}\cap\mathsf{SN}\subseteq
            \mathsf{SAT}$;
        \item if $X\in\mathsf{SAT}$, $t'\in X$ and
            $t\to_\mathrm{WHR}t'$, then $t\in X$.
    \end{itemize}

    Let there be the following relation
    $\mathcal{R}_\tau^\zeta\in\mathsf{SAT}$, where
    $\zeta$ is a mapping from type variables into
    $\mathsf{SAT}$:
    \begin{itemize}
        \item $t\in\mathcal{R}_\alpha^\zeta$ iff
            $t\in\zeta(\alpha)$;
        \item $t\in\mathcal{R}_\top^\zeta$ iff
            $t\in\mathsf{SN}$;
        \item $t\in\mathcal{R}_{T_1\times T_2}^\zeta$
            iff $t\in\mathsf{SN}$ and
            $\pi_i\,t\in\mathcal{R}_{T_i}^\zeta$
            for $i=1,2$;
        \item $t\in\mathcal{R}_{S\to T}^\zeta$ iff
            $t\in\mathsf{SN}$ and for all
            $s\in\mathcal{R}_S^\zeta$,
            $t\,s\in\mathcal{R}_T^\zeta$;
        \item $t\in\mathcal{R}_{\forall\alpha.T
            }^\zeta$ iff $t\in\mathsf{SN}$ and
            $t\in\mathcal{R}_T^{\zeta,\alpha\mapsto X}$
            for all $X\in\mathsf{SAT}$.
    \end{itemize}

    Indeed, $\mathcal{R}_\tau^\zeta\in\mathsf{SAT}$.
    Furthermore, the following property holds: if
    $t\in\mathcal{R}_{\forall\vec{\alpha}.T}^\zeta$,
    then $m\llbracket()\vdash\mathrm{nf}(t)\rrbracket
    \in\mathcal{R}^\zeta_{m\llbracket[\vec{S}/
    \vec{\alpha}]\vdash T\rrbracket}$ for all
    $m$ s.t. all terms inside it are in corresponding
    $\mathcal{R}_\forall^\zeta$.

    Then we can prove our main lemma: if
    $\Gamma\vdash t:T$ and $\gamma$ is a substitution
    s.t. $\gamma\,x\in\mathcal{R}_{\Gamma(x)}^\zeta$ for
    all $x\in\Gamma$, then $\gamma\,t\in
    \mathcal{R}_T^\zeta$.
\end{proof}

\begin{theorem}[Conservativity over STLC]
    Normal forms of closed terms typable in $\SSTLC$ are
    typable in STLC.
\end{theorem}
\begin{proof}
    Thanks to type preservation, normal form
    $\mathrm{nf}(t)$ is typable in $\SSTLC$. By
    definition, it is a value; therefore, it has no
    reinterpretation subterms and let-bindings, hence is
    typable in STLC.
\end{proof}

\begin{theorem}
    Global type inference for $\SSTLC$ is decidable.
\end{theorem}
\begin{proof}
    It is enough to extend any unification-based type
    inference algorithm for the Hindley-Milner type
    system to support models and reinterpretation.
    \begin{itemize}
        \item Since type operators in models are
            provided explicitly, it is enough to
            \textit{check} that their fields are typable
            with expected types.
        \item Since polymorphic functions do not depend
            on the local context, their type is fully
            resolved when they are added to the context.
        \item For case $m\llbracket x\rrbracket$, and
            $(x:\forall\vec{\alpha}.T)\in\Gamma$, we can
            directly compute $m\llbracket[\vec{\eta}/
            \vec{\alpha}]\vdash T\rrbracket$ (since
            $m$ is known and $T$ is fully resolved)
            where $\vec{\eta}$ are fresh metavariables.
    \end{itemize}
\end{proof}

Thus, in this section, we have introduced a sound type
system based on the Hindley-Milner type system which is
capable of interpreting its own closed terms in an
arbitrary ``cartesian closed category'' definable inside
it. This system is not truly \textit{reflective} since
models are not fully internal; they require a special
form of type operators to work, which is not accessible
to the rest of the type system. To internalize model
definition, we move on to a dependently typed system in
the next section.

\section{Almost reflective dependently typed language:
    \texorpdfstring{$\twoLTT$}{[[2]]LTT}}
\label{twoLTT}

Now we are ready to take on the main challenge: design
a minimal dependently typed reflective language, and
provide internal models of itself to yield universe
polymorphism, external parametricity and do synthetic
mathematics inside it. As a starting point, we take
Martin-L\"of type theory with $\Pi-$,$\Sigma-$,
$\mathrm{Id}-$ types, natural number type and infinite
non-cumulative hierarchy of universes. As discussed in
Section~\ref{no-go-section}, the best practical
solution to reinterpretation here would be to forbid
reductions under it and reinterpret terms as-is, and
constrain models using $\mathrm{Id}$-types. Still, in
comparison to non-dependently typed case, there are
infinitely more problems to solve if we want to define
a truly reflective system:

\begin{enumerate}
    \item If we na\"{\i}vely adopt an approach to
        typing of models from Section~\ref{SSTLC}, in
        order to specify a model, we would need to
        specify countably many models, one per each
        universe. And even if we did, to which universe
        does the type of models belong?
    \item Now that we have computations in types, we
        have to deal with conversion rule under
        reinterpretation:
        \[
            \dfrac{t:T\quad T\equiv T'}{t:T'}
            \stackrel{?}{\leadsto}
            \dfrac{\llbracket t\rrbracket:\llbracket T
            \rrbracket\quad\llbracket T\rrbracket\equiv
            \llbracket T'\rrbracket}{\llbracket t
            \rrbracket:\llbracket T'\rrbracket}
        \]
        The problem is that, after interpretation,
        $\llbracket T \rrbracket$ is only
        \textit{propositionally} equal to $\llbracket
        T'\rrbracket$.
    \item If we want a truly reflective system, we need
        to know what to do with
        \textit{reinterpretation of reinterpretation}:
        \[
            \lambda x.\llbracket\llbracket x
            \rrbracket_1\rrbracket_2
            \quad\text{and}\quad
            \llbracket\lambda x.\llbracket x
            \rrbracket_1\rrbracket_2
        \]
        In the first case, we perform two subsequent
        interpretations on a yet unknown $x$; in the
        second case, we interpret the process of
        interpretation itself. What should this mean?
\end{enumerate}

To deal with universe sizing issues, consider the
following observation: in a system without level
binders, each term uses only a finite number of levels.
Hence to reinterpret this term, we need to only provide
a finite number of universe models, which can together
fit into some larger universe of a finite level. No
universe polymorphism required! However, we now have to
track the maximal level which appears in a term in our
typing judgements.

To interpret conversion rule, we start by first
pinpointing the locations where conversion rule is
actually used; in order to do this ergonomically, we
use bidirectional type checking. At the moment of
conversion, interpretation has enough information to
know both $T$ and $T'$; thus we can translate a chain
of reductions performed during definitional equality
check into a chain of propositional equalities taken
from coherence conditions proven in the definition of a
model. Then, a point of conversion is wrapped in
equality elimination with synthesized witness as a
scrutinee.

In reinterpreting reinterpretation, we actually only
deal with ``double reinterpretation'': cases like
$\lambda x.\llbracket\llbracket x\rrbracket_1
\rrbracket_2$ could be dealt with by \textit{merging}
two models into a new one, a technique reminiscent of
\textit{Collapsing towers of interpreters}
\cite{Amin2017}; it is not required in our system since
the innermost reinterpretation is allowed to compute in
this case. As for interpreting interpretation itself as
in $\llbracket\lambda x.\llbracket x \rrbracket_1
\rrbracket_2$, we explicitly forbid such constructions
by introducing a family of modalities $\Box_\kappa$
meaning ``no interpretation was used in construction of
this term, and all levels appearing in it are not above
$\kappa$''. Terms have to have a boxed type to be
reinterpreted. In the notation for modality and general
typing discipline for it, we largely follow Gratzer et
al \cite{Gratzer2019} with a few important differences
discussed later. We also use $\Box$ to relax a ban on
computations inside reinterpretation by allowing
``splicing'' boxed expressions into other boxed
expressions; inside splices, computation is allowed to
bring terms to the boxed form for further
interpretation. Such a typing discipline brings us very
close to two-level type theories [CITE], but we found
notation via modalities to be more convenient.

\begin{figure}
\centering
\begin{bnf}
$\Gamma$ : Context ::=
    | $\cdot$ : empty context
    | $\Gamma,x\,\colon T$ : variable extension
    | $\Gamma.\,{}_\ell\text{\faLock}$
    : context locking
    ;;
$T$, $S$ : Type ::=
    | $\mathcal{U}_\ell$ : universe
    | $\mathrm{El}(t)$ : universal type family
    | $\mathrm{Nat}_\ell$ : natural number type
    | $(x\,\colon S)\times T$ : $\Sigma$-type
    | $(x\,\colon S)\to T$ : $\Pi$-type
    | $t_1=_T t_2$ : identity type
    | $\Box_\ell\,T$ : type of interpretable terms
    | $t\,\llbracket T\rrbracket$
        : interpretation of types
    ;;
$\ell$, $\kappa$, $j$ : Level :in: $\mathbb{N}$
    : universe levels ;;
$t$, $s$ : Term ::=
    | $x$ // $y$ // $z$ // $\ldots$
        : variables
    | $\mathtt{U}_\ell$ : code of a universe
    | $\mathtt{Nat}_\ell$
        : code of a natural number type
    | $0$ : zero
    | $\mathrm{S}\,t$ : successor
    | $\mathrm{natind}_{x.T}\,t\,\{s_0,s_S\}$
        : induction
    | $\mathtt{Sigma}(x\,\colon s)t$
        : code of a $\Sigma$-type
    | $(s, t)$ : pairing
    | $\pi_1\,t$ // $\pi_2\,t$ : projections
    | $\mathtt{Pi}(x\,\colon s)t$
        : code of a $\Pi$-type
    | $\lambda x.\,t$ : abstraction
    | $t\,s$ : application
    | $\mathtt{Id}\,t\,t_1\,t_2$
        : code of an identity type
    | $\mathrm{refl}$ : reflexivity
    | $\mathrm{J}_{xy.S}\,t\,s$
        : equality elimination
    | $\mathtt{Box}_\ell\,t$ : code of a box type
    | $[t]_\text{\faLock}$ : locking
    | $[t]_\text{\faUnlock}$ : unlocking
    | $t\,\llbracket s\rrbracket$
        : interpretation of terms
    | $t\,\colon T$ : ascription
    ;;
\end{bnf}
\caption{$\twoLTT$ grammar}
\label{twoltt-grammar}
\end{figure}

Grammar rules given in Figure~\ref{twoltt-grammar}
reflect discussion above. Note that there is no
separate type of models $\mathcal{M}$: instead of there
being a single builtin type, there is a family of
model types $\mathcal{M}_{\ell_0,\ldots,\ell_\kappa}^{
    \ell'_0,\ldots,\ell'_\kappa}$ where each concrete
type is \emph{definable} in $\twoLTT$; types in a
family are indexed by a bunch of universe levels whose
exact role will be explained later; a general idea is
that $\ell_j$ is a level which types in $\mathcal{U}_j$
would have after interpretation in such a model.

Typing rules are also a straightforward adaptation of a
general recipe for bidirectional typing of
dependently-typed systems \cite{Thiago2025} and are
given in Figures~\ref{twoltt-wf}-\ref{twoltt-typing}.
Important departure from standard bidirectional
recipe ``constructors are checked, eliminators are
inferred'' is that levels of types are always inferred
since we do not have cumulativity of universes and do
not have a lifting operator, thus we cannot check
levels of $\Sigma$ and $\Pi$. Also note that exact
model type is always inferred, both in interpretation
of types and in interpretation of terms: even though
$\mathcal{M}$ would be some huge nested $\Sigma$-type,
it is still more correct to \textit{match} on it
instead of asking to provide levels explicitly.

As already said, reinterpretation is constrained to act
on boxed types; in addition, we require a non-locked
context for reinterpretation. Together these
constraints guarantee that, indeed, boxed terms do not
contain reinterpretation subterms. An important
departure from $\mathrm{MLTT}_\text{\faLock}$
\cite{Gratzer2019}, aside from adding a family of
modalities, lies in the way context is locked in
$[\_]_\text{\faLock}$ and unlocked in
$[\_]_\text{\faUnlock}$: while locking, we require the
outer context to be unlocked; while unlocking, instead
of erasing all locks, we erase the context up to the
first lock. Effectively, this enforces a constraint
that typing context never contains more than a single
lock. Thus, for example, our $\Box_\ell$ is not a
comonad, but axiom K from modal logic still
holds\footnote{Interestingly enough, a similar
K-modality is first described by Birkedal and others in
``Modal dependent type theory and dependent right
adjoints'' \cite{Birkedal2020}. Coincidentally, they
refine typing judgements with universe levels, too!
They have a single modality, though.}:
\[
    \dfrac{\cdot\vdash_\ell A:\ell_A
    \quad\cdot\vdash_\ell B:\ell_B}{
        \cdot\vdash_\kappa\lambda f x.[[f
        ]_\text{\faUnlock}[x]_\text{\faUnlock}
        ]_\text{\faLock}:\Box_\ell(A\to B)\to
            \Box_\ell A\to\Box_\ell B
    }(\ell_A\sqcup\ell_B\leqslant\ell)
\]
In this way, unlocking truly gains the semantics of
``splicing'' in the sense that it can only contain
boxed values from the ``outer world''.

\begin{figure}
\centering
\begin{mathpar}
    \inferrule[WF-Nil]{}{\cdot\vdash_\ell}
    \and
    \inferrule[WF-Var]{
        \Gamma\vdash_\ell\\
        \Gamma\vdash_\ell T:\kappa
    }{\Gamma,x:T\vdash_\ell}(\kappa\leqslant\ell)
    \and
    \inferrule[WF-Lock]{
        \Gamma\vdash_\ell\\
        \text{\faLock}\not\in\Gamma
    }{\Gamma.\,{}_\ell\text{\faLock}\vdash_\kappa}
    \and
    \inferrule[Level-El]
        {\Gamma\vdash_\ell t\Rightarrow
            \mathcal{U}_\kappa}
        {\Gamma\vdash_\ell\mathrm{El}(t):\kappa}
        (\kappa+1\leqslant\ell)
    \and
    \inferrule[Level-Uni]{\Gamma\vdash_\ell}
        {\Gamma\vdash_\ell\mathcal{U}_\kappa:\kappa+1}
    \and
    \inferrule[Level-Nat]{\Gamma\vdash_\ell}
        {\Gamma\vdash_\ell\mathrm{Nat}_\kappa:\kappa}
    \and
    \inferrule[Level-Sigma]{
        \Gamma\vdash_\ell S:j \\
        \Gamma,x:S\vdash_\ell T:\kappa
    }{\Gamma\vdash_\ell(x:S)\times T:j\sqcup\kappa}
    \and
    \inferrule[Level-Pi]{
        \Gamma\vdash_\ell S:j \\
        \Gamma,x:S\vdash_\ell T:\kappa
    }{\Gamma\vdash_\ell(x:S)\to T:j\sqcup\kappa}
    \and
    \inferrule[Level-Id]{
        \Gamma\vdash_\ell T:\kappa \\
        \Gamma\vdash_\ell t_1\Leftarrow T \\
        \Gamma\vdash_\ell t_2\Leftarrow T
    }{\Gamma\vdash_\ell t_1=_T t_2:\kappa}
    \and
    \inferrule[Level-Box]{
        \text{\faLock}\not\in\Gamma\\
        \Gamma.\,{}_\ell\text{\faLock}\vdash_\kappa T:j
    }{\Gamma\vdash_\ell\Box_\kappa\,T:\kappa}
    (j\leqslant\kappa)
    \and
    \inferrule[Level-TyInterp]{
        \Gamma\vdash_\ell t\Rightarrow\mathcal{M}_{
            \ell_0,\ldots,\ell_\kappa}^{
            \ell'_0,\ldots,\ell'_\kappa}\\
        \text{\faLock}\not\in\Gamma\\
        \Gamma.{}_\ell\text{\faLock}\vdash_\kappa T:j
    }{\Gamma\vdash_\ell t\llbracket T\rrbracket:\ell_j}
    (j\leqslant\kappa)
\end{mathpar}
\caption{
    Context well-formedness \& level placement rules}
\label{twoltt-wf}
\end{figure}

\begin{figure}
\centering
\begin{mathpar}
    \inferrule[Infer-Uni]{\Gamma\vdash_\ell}{
        \Gamma\vdash_\ell\mathtt{U}_\kappa\Rightarrow
            \mathcal{U}_{\kappa+1}
    }
    \and
    \inferrule[Infer-Nat]{\Gamma\vdash_\ell}{
        \Gamma\vdash_\ell\mathtt{Nat}_\kappa\Rightarrow
            \mathcal{U}_\kappa}
    \and
    \inferrule[Check-Zero]{\Gamma\vdash_\ell}{
        \Gamma\vdash_\ell 0\Leftarrow
            \mathrm{Nat}_\kappa
    }
    \and
    \inferrule[Check-Succ]
        {\Gamma\vdash_\ell t\Leftarrow
            \mathrm{Nat}_\kappa}
        {\Gamma\vdash_\ell\mathrm{S}\,t\Leftarrow
            \mathrm{Nat}_\kappa}
    \and
    \inferrule[Infer-NatInd]{
        \Gamma\vdash_\ell t\Rightarrow
            \mathrm{Nat}_\kappa\\
        \Gamma,x:\mathrm{Nat}_\kappa\vdash_\ell T:j\\
        \Gamma\vdash_\ell s_0\Leftarrow
            T[(0:\mathrm{Nat}_\kappa)/x]\\
        \Gamma\vdash_\ell s_S\Leftarrow
                (y:\mathrm{Nat}_\kappa)\to
                (z:T[y/x])\to
                T[(\mathrm{S}\,y:\mathrm{Nat}_\kappa)
                /x]
    }{\Gamma\vdash_\ell
        \mathrm{natind}_{x.T}\,t\,\{s_0,s_S\}
        \Rightarrow T[t/x]
    }(j\leqslant\ell)
    \and
    \inferrule[Infer-Sigma]{
        \Gamma\vdash_\ell s\Rightarrow
            \mathcal{U}_j\\
        \Gamma,x:\mathrm{El}(s)\vdash_\ell t
            \Rightarrow\mathcal{U}_\kappa
    }{\Gamma\vdash_\ell\mathtt{Sigma}(x:s)t\Rightarrow
        \mathcal{U}_{j\sqcup\kappa}}
    \and
    \inferrule[Check-Pair]{
        \Gamma\vdash_\ell s\Leftarrow S\\
        \Gamma\vdash_\ell t\Leftarrow T[(s:S)/x]
    }{\Gamma\vdash_\ell(s,t)\Leftarrow(x:S)\times T}
    \and
    \inferrule[Infer-$\mathrm{Proj}_1$]
        {\Gamma\vdash_\ell t\Rightarrow(x:S)\times T}
        {\Gamma\vdash_\ell \pi_1\,t\Rightarrow S}
    \and
    \inferrule[Infer-$\mathrm{Proj}_2$]
        {\Gamma\vdash_\ell t\Rightarrow(x:S)\times T}
        {\Gamma\vdash_\ell \pi_2\,t\Rightarrow
            T[(\pi_1\,t)/x]}
    \and
    \inferrule[Infer-Pi]{
        \Gamma\vdash_\ell s\Rightarrow
            \mathcal{U}_j\\
        \Gamma,x:\mathrm{El}(s)\vdash_\ell t
            \Rightarrow\mathcal{U}_\kappa
    }{\Gamma\vdash_\ell\mathtt{Pi}(x:s)t\Rightarrow
        \mathcal{U}_{j\sqcup\kappa}}
    \and
    \inferrule[Check-Abs]
        {\Gamma,x:S\vdash_\ell t\Leftarrow T}
        {\Gamma\vdash_\ell\lambda x.\,t\Leftarrow
            (x:S)\to T}
    \and
    \inferrule[Infer-App]{
        \Gamma\vdash_\ell t\Rightarrow (x:S)\to T\\
        \Gamma\vdash_\ell s\Leftarrow S
    }{\Gamma\vdash_\ell t\,s\Rightarrow
        T[(s:S)/x]}
    \and
    \inferrule[Infer-Id]{
        \Gamma\vdash_\ell t\Rightarrow
            \mathcal{U}_\kappa\\
        \Gamma\vdash_\ell t_1\Leftarrow\mathrm{El}(t)
        \\\Gamma\vdash_\ell t_2\Leftarrow\mathrm{El}(t)
    }{\Gamma\vdash_\ell\mathtt{Id}\,t\,t_1\,t_2
        \Rightarrow\mathcal{U}_\kappa}
    \and
    \inferrule[Check-Refl]{t_1\equiv t_2}
        {\Gamma\vdash_\ell\mathrm{refl}\Leftarrow
            t_1=_T t_2}
    \and
    \inferrule[Infer-J]{
        \Gamma\vdash_\ell t\Rightarrow t_1=_T t_2\\
        \Gamma,x:T,y:t_1=_T x\vdash_\ell S:\kappa\\\\
        \Gamma\vdash_\ell s\Leftarrow
            S[(t_1:T)/x,
                (\mathrm{refl}:t_1=_T t_1)/y]
    }{\Gamma\vdash_\ell\mathrm{J}_{xy.S}\,t\,s
        \Rightarrow S[(t_2:T)/x, t/y]
    }(\kappa\leqslant\ell)
    \and
    \inferrule[Infer-Var]{
        \text{\faLock}\not\in\Delta\\
        \Gamma,x:T,\Delta\vdash_\ell
    }{\Gamma,x:T,\Delta\vdash_\ell x\Rightarrow T}
    \and
    \inferrule[Infer-Box]{
        \text{\faLock}\not\in\Gamma\\
        \Gamma.\,{}_\ell\text{\faLock}\vdash_\kappa t
            \Rightarrow\mathcal{U}_j
    }{\Gamma\vdash_\ell\mathtt{Box}_\kappa\,t
        \Rightarrow\mathcal{U}_\kappa}
    (j\leqslant\kappa)
    \and
    \inferrule[Check-Lock]{
        \text{\faLock}\not\in\Gamma\\
        \Gamma.\,{}_\ell\text{\faLock}\vdash_\kappa t
            \Leftarrow T
    }{\Gamma\vdash_\ell[t]_\text{\faLock}\Leftarrow
        \Box_\kappa T}
    \and
    \inferrule[Infer-Unlock]{
        \text{\faLock}\not\in\Delta\\
        \Gamma\vdash_\kappa t\Rightarrow\Box_j T
    }{\Gamma.{}_\kappa\text{\faLock},\Delta
        \vdash_\ell[t]_\text{\faUnlock}\Rightarrow T
    }(j\leqslant\ell)
    \and
    \inferrule[Infer-Interp]{
        \text{\faLock}\not\in\Gamma\\
        \Gamma\vdash_\ell s\Rightarrow\mathcal{M}_{
            \ell_0,\ldots,\ell_k}^{
            \ell'_0,\ldots,\ell'_k}\\
        \Gamma\vdash_\ell t\Rightarrow \Box_\kappa T
    }{\Gamma\vdash_\ell s \llbracket t\rrbracket
    \Rightarrow s \llbracket T\rrbracket
    }
    \and
    \inferrule[Infer-Ascr]{
        \Gamma\vdash_\ell T:\kappa\\
        \Gamma\vdash_\ell t\Leftarrow T
    }{\Gamma\vdash_\ell t:T\Rightarrow T}
    (\kappa\leqslant\ell)
    \and
    \inferrule[Check-Inferred]{
        \Gamma\vdash_\ell t\Rightarrow S\\ S\equiv T
    }{\Gamma\vdash_\ell t\Leftarrow T}
\end{mathpar}
\caption{Bidirectional term typing rules}
\label{twoltt-typing}
\end{figure}

\subsection{Computation rules and models of dependent
type theories}

Note that discussion above was ambient to the choice of
formalism to present models. Now it is finally time to
fix the formalism since we're going to need it to
define the type of models $\mathcal{M}$.

In line with current state of the art, we use
Categories-with-Families \cite{castellan2020cwf} (or
CwFs for short) as it is currently one of the most
convenient frameworks to specify models of dependent
type theories. Here is a self-contatined definition of
a category with families, given for clear reference in
subsequent internalization:

\begin{definition}[Category]
A category is the following datum\footnote{Of course,
cartesian closed category (\ref{ccc}) is also a
category. We repeat parts of definitions for
self-containment reasons.}:
\begin{enumerate}
    \item A class of \textit{objects} $\mathrm{Ob}$ and
        a family of \textit{morphisms}
        $\mathrm{Hom}(S, T)$ indexed by
        $S,T\in\mathrm{Ob}$, where $f\in\mathrm{Hom}(S,
        T)$ is denoted by $f:S\to T$. In addition, $S$
        is called \textit{source} of $f$, while $T$ is
        \textit{target} of $f$.
    \item For each $A\in\mathrm{Ob}$ there's an
        \textit{identity morphism} $\mathrm{id}_A:
        A\to A$; for each $f:A\to B$ and $g:B\to C$,
        there's their \textit{composition} $g\circ
        f:A\to C$. If a source-target of an identity
        morphism is clear from the context, the
        subscript is elided. Also, the following
        properties hold:
        \begin{align*}
            \mathrm{id}\circ f = f = f\circ\mathrm{id}
            & &
            (f \circ g) \circ h = f \circ (g \circ h)
        \end{align*}
        whenever all compositions are well-defined.
\end{enumerate}
\end{definition}

\begin{definition}[Terminal object]
    Terminal object of category $\mathcal{C}$ is an
    object $\top\in\mathrm{Ob}_\mathcal{C}$ such that
    for each $A\in\mathrm{Ob}_\mathcal{C}$ there's a
    \textbf{unique} morphism $!_A:A\to\top$.
\end{definition}

\begin{definition}[Presheaf]
    A presheaf $P$ on a category $\mathcal{C}$ (written
    as $P:\mathcal{C}^\mathrm{op}\to\mathrm{Set}$) is a
    mapping which takes objects of $\mathcal{C}$ to
    sets, and morphisms of $\mathcal{C}$ to functions
    between sets, so that if $f:A\to_\mathcal{C} B$,
    then $P(f)$ is a function from $P(B)$ to $P(A)$.
    Also, the following properties hold:
    \begin{align*}
        P(\mathrm{id}_A)=\mathrm{id}_{P(A)} & &
        P(g\circ f)=P(f)\circ P(g)
    \end{align*}
    Where $\mathrm{id}_{P(A)}$ is an identity function
    of a set $P(A)$ and $P(f)\circ P(g)$ is a
    composition of functions $P(f)$ and $P(g)$.
\end{definition}

\begin{definition}[Category of elements]
    Given a presheaf $P:\mathcal{C}^\mathrm{op}\to
    \mathrm{Set}$, a category of elements
    $\int_\mathcal{C}P$ is a category where objects are
    pairs $(A, x)$ such that
    $A\in\mathrm{Ob}_\mathcal{C}$ and $x\in P(A)$, and
    morphisms from $(A, x)$ to $(B, y)$ are morphisms
    $f:A\to_\mathcal{C} B$ such that $P(f)(y)=x$.
\end{definition}

\begin{definition}[Category with families]
    \label{cwfdef}
    A category with families is the following datum:
    \begin{enumerate}
        \item A \textit{category of contexts}
            $\mathcal{C}$ which is a category with a
            terminal object;
        \item A \textit{presheaf of types}
            $\mathrm{Ty}:\mathcal{C}^\mathrm{op}\to
            \mathrm{Set}$. Given a morphism $\gamma:
            \Gamma\to_\mathcal{C}\Delta$ and a
            \textit{type} $T\in\mathrm{Ty}(\Delta)$,
            element $\mathrm{Ty}(\gamma)(T)\in
            \mathrm{Ty}(\Gamma)$ is usually written as
            $T[\gamma]$ for brevity.
        \item A \textit{presheaf of terms}
            $\mathrm{Tm}:\left(\int_\mathcal{C}
            \mathrm{Ty}\right)^\mathrm{op}\to
            \mathrm{Set}$. Same as with types,
            $\mathrm{Tm}(\gamma)(t)\in\mathrm{Tm}(
            \Gamma, T)$ is usually written as
            $t[\gamma]$ for brevity.
        \item For each
            $\Gamma\in\mathrm{Ob}_\mathcal{C}$ and
            $T\in\mathrm{Ty}(\Gamma)$, there is an
            \textit{extended context} $\Gamma.T\in
            \mathrm{Ob}_\mathcal{C}$ together with a
            \textit{weakening morphism} $p:\Gamma.T\to
            \Gamma$ and \textit{variable element}
            $q:\mathrm{Tm}(\Gamma.T, T[p])$. In
            addition, for all $\Delta$, $\gamma:\Delta
            \to\Gamma$ and $t\in\mathrm{Tm}(\Delta,
            T[\gamma])$, there is a unique
            \textit{extended substitution} $\langle
            \gamma,t\rangle:\Delta\to\Gamma.T$ such
            that the following holds:
            \begin{align*}
                p\circ\langle\gamma,t\rangle=\gamma & &
                q[\langle \gamma, t \rangle] = t
            \end{align*}
\end{enumerate}
\end{definition}

A definition of a Cwf (def. \ref{cwfdef}) clearly
captures the bare-bones structure of a dependently
typed theory: objects of $\mathcal{C}$ represent
contexts, while morphisms represent substitutions
mapping one context to another. Substitution also acts
both on types and terms; given a correct type $T$ in
context $\Gamma$, one can extend a context to gain new
context $\Gamma,x:T$ in which there's a well-typed term
$x:T$; there's a way to obtain substitutions into an
extended context by providing a term suitable for $x$.

This base definition can be readily extended to
incorporate type (and term) formers of any system. A
CwF for $\twoLTT$, then, has to include $\mathrm{Ty}$-
and $\mathrm{Tm}$- formers for analogues of universes,
$\Pi$- and $\Sigma$- types, natural numbers and
identity types. Interpretation for $\Box_\ell$ is not
needed since we succesfully prevented boxed types from
appearing inside boxed terms via rigid \faLock-ing
discipline\footnote{That's also the reason we dubbed
this system \textit{almost} reflective, since its
internal reinterpretation doesn't cover modalities and
reinterpretation itself.}. Thus, contexts which need to
be interpreted are constructed via vanilla context
extensions from empty context, and a usual CwF is
enough for our purposes. Since actual internalized
definition of our CwF largely repeats typing and
computation rules from
Figures~\ref{twoltt-wf}-\ref{twoltt-typing} and
Figures~\ref{twoltt-typcomp}-\ref{twoltt-comp} in
internal language, we do not provide it here; to get a
bit more details, you can refer to our
Appendix~\ref{twoltt-internal-model-def}.

Finally, on to computation rules as given in
Figures~\ref{twoltt-typcomp}-\ref{twoltt-comp}.
Mostly they are standard MLTT computation rules for
bidirectional typing (without $\eta$-equality for now).
Importantly enough, we \textbf{do not compute inside
locked terms and inside boxed types} in order to not
stumble into a variant of our no-go theorem \ref{no-go}
again (modal variant is spelled out in discussion, see
Section~\ref{global-elements}).

\begin{figure}
\centering
\begin{mathpar}
    \inferrule{t\leadsto t'}
    {\mathrm{El}(t)\leadsto\mathrm{El}(t')}
    \and
    \inferrule{}{\mathrm{El}(\mathtt{U}_\kappa)\leadsto
        \mathcal{U}_\kappa}
    \and
    \inferrule{}{\mathrm{El}(\mathtt{Nat}_\kappa)
        \leadsto\mathrm{Nat}_\kappa}
    \and
    \inferrule{}{\mathrm{El}(\mathtt{Sigma}(x:s)t)
        \leadsto(x:\mathrm{El}(s))\times\mathrm{El}(t)}
    \and
    \inferrule{}{\mathrm{El}(\mathtt{Pi}(x:s)t)\leadsto
            (x:\mathrm{El}(s))\to\mathrm{El}(t)}
    \and
    \inferrule{}{\mathrm{El}(\mathtt{Id}\,t\,t_1\,t_2)  
        \leadsto t_1=_{\mathrm{El}(t)}t_2}
    \and
    \inferrule{}{\mathrm{El}(\mathtt{Box}_\kappa t)
        \leadsto\Box_\kappa\mathrm{El}(t)}
    \and
    \inferrule{S\leadsto S'}
        {(x:S)\times T\leadsto (x:S')\times T}
    \and
    \inferrule{T\leadsto T'}
    {(x:S)\times T\leadsto(x:S)\times T'}
    \and
    \inferrule{S\leadsto S'}
    {(x:S)\to T\leadsto (x:S')\to T}
    \and
    \inferrule{T\leadsto T'}
    {(x:S)\to T\leadsto(x:S)\to T'}
    \and
    \inferrule{T\leadsto T'}
        {t_1=_T t_2\leadsto t_1=_{T'} t_2}
    \and
    \inferrule{t_1\leadsto t_1'}
        {t_1=_T t_2\leadsto t_1'=_T t_2}
    \and
    \inferrule{t_2\leadsto t_2'}
        {t_1=_T t_2\leadsto t_1=_T t_2'}
    \and
    \inferrule{}{
        t\llbracket T\rrbracket\leadsto
        t.\mathrm{Tm}(t.\cdot, t\llbracket\cdot\vdash
        T\rrbracket)}
\end{mathpar}
    \caption{$\twoLTT$ typal computation rules}
\label{twoltt-typcomp}
\end{figure}

\begin{figure}
\centering
\begin{mathpar}
    \inferrule{t\leadsto t'}{S\,t\leadsto S\,t'}
    \and
    \inferrule{t\leadsto t'}{
        \mathrm{natind}_{x.T}t\{s_0,s_S\}\leadsto
        \mathrm{natind}_{x.T}t'\{s_0,s_S\}
    }
    \and
    \inferrule{}{\mathrm{natind}_{x.T}
        (0:\mathrm{Nat}_\ell)\{s_0,s_S\}
        \leadsto (s_0:T[(0:\mathrm{Nat}_\ell)/x]
        )}
    \and
    \inferrule{}{\mathrm{natind}_{x.T}
        (S\,t:\mathrm{Nat}_\ell)\{s_0,s_S\}
        \leadsto (s_S:(x:\mathrm{Nat}_\ell)\to T\to
            T[(S x)/x])\,t\,(\mathrm{natind
            }_{x.T}(t:\mathrm{Nat}_\ell)\{s_0,s_S\})}
    \and
    \inferrule{s\leadsto s'}{\mathtt{Sigma}(x:s)t
        \leadsto\mathtt{Sigma}(x:s')t}
    \and
    \inferrule
        {t\leadsto t'}
        {\mathtt{Sigma}(x:s)t\leadsto
            \mathtt{Sigma}(x:s)t'}
    \and
    \inferrule{s\leadsto s'}{(s,t)\leadsto(s',t)}
    \and
    \inferrule{t\leadsto t'}{(s,t)\leadsto(s,t')}
    \and
    \inferrule{t\leadsto t'}{\pi_i t\leadsto \pi_i t'}
    \and
    \inferrule{}
    {\pi_i((t_1,t_2):(x:S)\times T)\leadsto t_i}
    \and
    \inferrule{s\leadsto s'}{\mathtt{Pi}(x:s)t
        \leadsto\mathtt{Pi}(x:s')t}
    \and
    \inferrule
        {t\leadsto t'}
        {\mathtt{Pi}(x:s)t\leadsto
            \mathtt{Pi}(x:s)t'}
    \and
    \inferrule
        {t\leadsto t'}
        {\lambda x.t\leadsto\lambda x.t'}
    \and
    \inferrule{t\leadsto t'}{t\,s\leadsto t'\,s}
    \and
    \inferrule{s\leadsto s'}{t\,s\leadsto t\,s'}
    \and
    \inferrule{}
        {((\lambda x.t):(x:S)\to T)s\leadsto t[(s:S)/x]}
    \and
    \inferrule{t\leadsto t'}{\mathtt{Id}\,t\,t_1\,t_2
        \leadsto\mathtt{Id}\,t'\,t_1\,t_2}
    \and
    \inferrule{t_1\leadsto t_1'}{
        \mathtt{Id}\,t\,t_1\,t_2\leadsto
        \mathtt{Id}\,t\,t_1'\,t_2
    }
    \and
    \inferrule{t_2\leadsto t_2'}{
        \mathtt{Id}\,t\,t_1\,t_2\leadsto
        \mathtt{Id}\,t\,t_1\,t_2'
    }
    \and
    \inferrule{t\leadsto t'}{\mathrm{J}_{xy.S}\,t\,s
        \leadsto\mathrm{J}_{xy.S}\,t'\,s}
    \and
    \inferrule{}{\mathrm{J}_{xy.S}(\mathrm{refl}:
        t_1=_T t_2)\,s\leadsto(s:S[t_2/x,
            (\mathrm{refl}:t_1=_T t_2)/y])}
    \and
    \inferrule{t\leadsto t'}{(t:T)\leadsto(t':T)}
    \and
    \inferrule{t\leadsto t'}
    {[t]_\text{\faUnlock}\leadsto[t']_\text{\faUnlock}}
    \and
    \inferrule{}{[[t]_\text{\faLock}:\Box_\kappa T
    ]_\text{\faUnlock}\leadsto t:T}
    \and
    \inferrule{}{[[t]_\text{\faUnlock}]_\text{\faLock}
    \leadsto t}
    \and
    \inferrule{s\leadsto s'}{s\llbracket t\rrbracket
    \leadsto s'\llbracket t\rrbracket}
    \and
    \inferrule{t\leadsto t'}{s\llbracket t\rrbracket
    \leadsto s\llbracket t'\rrbracket}
    \and
    \inferrule{}{
        s\llbracket[t]_\text{\faLock}:\Box_\kappa T
        \rrbracket\leadsto
        s\llbracket\cdot\vdash t\Leftarrow T\rrbracket
    }
    \and
    \inferrule{ }{t:T\leadsto t}(t\text{ is
    synthesizing})
\end{mathpar}
\caption{$\twoLTT$ term computation rules}
\label{twoltt-comp}
\end{figure}

Also we need to give the rules for computing
reinterpretation; as it turns out,
\textit{reinterpreting terms as-is} allows to define
reinterpretation not as a redex, but rather as a couple
of total functions on types and terms, just as it was in
a non-dependently-typed case (Section~\ref{SSTLC}). To
be precise, we need two mutually recursive functions
for terms:
$m\llbracket\Gamma\vdash t\Leftarrow T\rrbracket$
computes on checkable terms, while $m\llbracket\Gamma
\vdash t\Rightarrow\rrbracket$ computes on synthesizing
terms and outputs an inferred type; both compute
elements of corresponding $\mathrm{Tm}(\Gamma,\_)$-s
from a CwF $m$. $m\llbracket\Gamma\vdash T\rrbracket$
computes an element of $\mathrm{Ty}(\Gamma)$. Another
auxiliary function is required to compute chain of
equalities to handle conversion rule; this function has
two subtle points:

\begin{itemize}
    \item At this point, we don't know if system is
        strongly normalizing yet, therefore we
        potentially might encounter an infinite chain of
        equalities.
    \item In systems where equality is proof-relevant
        (and we want to be able to develop in that
        direction), different reduction chains might
        lead to different equality witnesses, so such a
        function is potentially not well-defined.
\end{itemize}

Our solution is the following: let us fix a certain
evaluation strategy (for example, strict evaluation)
and employ it in normalization. Since terms we reduce
to find a witness are vanilla MLTT terms, they are
strongly normalizing, so witness extraction terminates.

TODO: provide interpretation definition.

\subsection{Metatheoretical properties of
\texorpdfstring{$\twoLTT$}{[[2]]LTT}}
\label{twoLTT-properties}

Once again, we only try to make sure of the most
essential properties: that computations are confluent,
that types during computations are preserved, all
computations terminate and that resulting type system
is a conservative extension of MLTT.

\begin{theorem}[Church-Rosser theorem for $\twoLTT$]
    Reflexive transitive closure of $(\leadsto)$ has
    Church-Rosser property.
\end{theorem}
\begin{proof}
    Classic proof via parallel reduction relation
    [CITE]. Since interpretation computes as a
    recursive functions on open terms, it commutes with
    substitution, which is crucial for confluence.

    TODO: write out parallel reduction?
\end{proof}

\begin{theorem}[Type preservation for $\twoLTT$]
    If $\Gamma\vdash t\Leftarrow T$ and $t\leadsto t'$,
    then $\Gamma\vdash t'\Leftarrow T$; analogously for
    $\Gamma\vdash t\Rightarrow T$.
\end{theorem}
\begin{proof}
    Note that, analogous to other dependently typed
    systems, a context reduction lemma holds which is
    crucial to handle reduction in domain of
    $\Pi$-types (where both $\Leftarrow$ and
    $\Rightarrow$ can stand in place of a colon):
    \[
        \dfrac{
            \Gamma,x:S,\Delta\vdash_\ell t:T
            \quad S\leadsto S'
        }{\Gamma,x:S',\Delta\vdash_\ell t:T}
    \]

    Then the proof proceeds by structural induction on
    terms and $(\leadsto)$. Note that, in a
    bidirectional system, one has to take care to
    structure rules in such a way that substitutions
    substitute terms that are inferred, so that
    inferrable terms stay inferrable and checkable
    terms stay checkable.

    An important difference from other dependently
    typed systems is, of course, interpretation. For
    this case, analogous to the proof of type
    preservation for $\SSTLC$
    (Section~\ref{SSTLC-properties}), a vital lemma
    holds:
    \[  \dfrac{
            \Gamma.{}_\ell\text{\faLock},\Delta
                \vdash_\kappa t\Leftarrow T
            \quad\Gamma\vdash_\ell s\Rightarrow
                \mathcal{M}_{\ell_0,\ldots,\ell_\kappa
                }^{\ell'_0,\ldots,\ell'_\kappa}
        }{\Gamma\vdash_\ell s\llbracket\Delta\vdash t
        \Leftarrow T\rrbracket\Rightarrow
            s.\mathrm{Tm}(s\llbracket\Delta\rrbracket,
                s\llbracket\Delta\vdash T\rrbracket)}
    \]
    Along with similar lemmas for
    $\_\llbracket\_\vdash\_\Rightarrow\rrbracket$ and
    $\_\llbracket\_\vdash\_\rrbracket$.
\end{proof}

\begin{theorem}[Identity model]
    For each $\ell$, there is an identity model
    $I_\ell:\mathcal{M}_{0,\ldots,\ell}^{1,\ldots,
    \ell+1}$ with a property that $I_\ell\llbracket
    [t]_\text{\faLock}:\Box_\kappa T\rrbracket\equiv
    t:T$ for each closed $t$.
\end{theorem}
\begin{proof}
    Take a universe at appropriate level as type of
    contexts; $\mathrm{Ty}_\ell(\Gamma)=
    \Gamma\to\mathcal{U}_\ell$;
    $\mathrm{Tm}_\ell(\Gamma, T)=(\gamma:\Gamma)\to
    T\,\gamma$. Define all of the remaining term
    interpretations to reinterpret to the same terms.
\end{proof}

\begin{theorem}[Strong normalization for $\twoLTT$]
    Closed terms typable in $\twoLTT$ are strongly
    normalizing with respect to $(\leadsto)$.
\end{theorem}
\begin{proof}
    We could reprove normalization from scratch, just
    how we did with $\SSTLC$. However, here we can
    provide a much simpler proof which delegates
    normalization to another type system. To be more
    precise, this system $\mathcal{S}$ is Martin-L\"of
    type theory with $\omega+\omega$ universes and
    first-class universe polymorphism allowed to range
    over first $\omega$ universes; we know it to be
    strongly normalizing [CITE]. Then, we define a
    translation function $[\_]$ which transforms types
    and terms of $\twoLTT$ into types and terms of
    $\mathcal{S}$: it preserves all vanilla grammar
    formers and only translates lock
    $[\_]_\text{\faLock}$ into a function from the type
    of models $\mathcal{M}_{\ell_0,\ldots,\ell_\kappa
    }^{\ell'_0,\ldots,\ell'_\kappa}$ into interpreted
    type. Since $\kappa$ is fixed, first-class universe
    polymorphism is enough to translate it into a
    function which accepts models for \emph{any} list
    of levels of fitting length. Then, unlocking
    $[\_]_\text{\faUnlock}$ translates into application
    of this function to identity model of fitting
    level; interpretation translates into application.

    Note that $t\leadsto t'$ implies
    $[t]\leadsto^+[t']$; this concludes the proof since
    any infinite chain of reductions in our system
    would imply an existence of infinite chain of
    reductions in $\mathcal{S}$.
\end{proof}

\begin{theorem}[Conservativity over MLTT]
    For MLTT context $\Gamma$, MLTT type $T$ and
    $\twoLTT$ term $t$, if
    $\Gamma\vdash t\Leftarrow T$, then
    $\Gamma\vdash_\mathrm{MLTT}\mathrm{nf}(t)\Leftarrow
    T$.
\end{theorem}
\begin{proof}
    Thanks to type preservation, normal form
    $\mathrm{nf}(t)$ is typable in $\twoLTT$. Since
    $\Gamma$ does not contain neither $\Box$-es nor
    \faLock-s, any subterm of type $\Box_\kappa T$
    has form $[\_]_\text{\faLock}$. Since $T$ is an MLTT
    type, every such locking is wrapped either in a
    $[\_]_\text{\faUnlock}$ or in a $\llbracket\_
    \rrbracket$. Second case is a redex and is
    impossible in a normal form; $[\_]_\text{\faUnlock}$
    is possible only in a locked context, so that means
    that there's another $[\_]_\text{\faLock}$ above it.
    Therefore $\mathrm{nf}(t)$ has no subterms of type
    $\Box_\kappa T$; then it has no unlocking and
    interpretation subterms and is an MLTT term and is
    typable in bidirectional MLTT (it is obtained from
    $\twoLTT$ by getting rid of rules involving
    modalities, locking and reinterpretation).
\end{proof}

\subsection{Useful reinterpretations inside
\texorpdfstring{$\twoLTT$}{[[2]]LTT}}
\label{twoLTT-applications}

Here we list some of the language features which
reinterpretation enables without any additional language
support.

\subsubsection{Crude, but effective stratification}
\label{twoLTT-displacement}

Due to McBride \cite{McBride:Crude2011,
McBride:Crude2012}, a technique of crude, but effective
stratification works as follows: instead of providing
complete support for level polymorphism, let's consider
a \textit{displacement operator} which, given a term $t$
using universes $U_i,U_j,U_k,\ldots$, produces a term
which is equal to $t$ except for universes which are
replaced with $U_{i+1},U_{j+1},U_{k+1},\ldots$. The
displaced term is also well-typed, however it is defined
<<one level higher>>. While not directly definable as a
single operator inside $\twoLTT$, for each
$\kappa,\ell$, one can define displacement operators
$\uparrow_\ell^\kappa:\mathcal{M}_{\kappa+0,
\kappa+1,\ldots,\kappa+\ell}^{\kappa+1,\kappa+2,\ldots,
\kappa+\ell+1}$ that shifts a boxed definition which
uses universes up to $\mathcal{U}_\ell$ by $\kappa$
levels up. In this model, category of contexts basic
presheaves are shifted by $\kappa$ in comparison to an
identity model; all other term interpretations are
identical to the ones given in an identity model.

A bit more details on definition of
$\uparrow_\ell^\kappa$ are given in
Appendix~\ref{twoltt-displacement-model-def}.

\subsubsection{Internal external parametricity}
\label{twoLTT-parametricity}

Demonstrated by Reynolds \cite{reynolds1983types},
parametricity is a property of a type system that, in
simple words, <<polymorphic function cannot
meaningfully depend on generic inputs>>. Thus any
polymorphic function of type $\sigma$ satisfies a
certain <<free theorem>> $\llbracket \sigma
\rrbracket$ defined via relational interpretation of a
type.

In recent works \cite{Bernardy:Parametricity2010,
Altenkirch:Parametricity2024}, a streamlined definition
is given for parametricity interpretation of
dependently-typed systems dubbed <<external
parametricity>>. Once again, this relational
interpretation is definable in our system
\textit{internally} with the help of a relational model
$\mathrm{Para}_\ell:\mathcal{M}_{1,\ldots,\ell+1}^{2,
\ldots,\ell+2}$. The shift by one level happens because
type $T:\mathcal{U}_\ell$ in such a model is
interpreted by a triple $(T_1\,T_2:\mathcal{U}_\ell)
\times(T_1\to T_2\to\mathcal{U}_\ell)$. Since
interpreted functions do not depend on local context,
we sidestep issues with relational context highlighted
by Altenkirch and others
\cite{Altenkirch:Parametricity2024}, thus result of the
interpretation is a proper proof of a free theorem.

A bit more details on definition of
$\mathrm{Para}_\ell$ are given in
Appendix~\ref{twoltt-parametricity-model-def}.

\subsubsection{Synthetic mathematics}
\label{twoLTT-synthetic}

As thoroughly discussed in a recent paper by Nawrocki
and others \cite{Nawrocki:Lean2026}, thanks to MLTT
being an internal language of certain
$(\infty,1)$-categories, we can do \textit{synthetic
mathematics} inside it: prove a theorem inside MLTT
about types; interpret it in a model of MLTT and yield
a proof about objects in that model. It is easy to see
that reinterpretation operator we have in $\twoLTT$ is
sufficient to do synthetic mathematics in it and
immediately apply results interpreted in a suitable
model.

\section{Discussion \& related work}

\subsection{What was done before us}

Original work by Lambek \cite{Lambek1980} sets the tone
of the future research, establishing correspondence
between type systems based on lambda calculus with
various kinds of categories, starting with
correspondence between simply typed lambda calculus and
cartesian closed categories.

Conal Elliott's ``Compiling to Categories''
\cite{Elliott:Compcat2017} shows how to interpret (some
subset of) Haskell in arbitrary \texttt{Category}
(again, definable in Haskell) via GHC compiler plugin.
That paper is a direct inspiration for this one. There
is an important difference, though: reinterpretation is
internal to our type systems, unlocking new patterns in
the usage of reinterpretation, including (but not
limited to) universe polymorphism, external
parametricity and more. However, Conal's approach is
more flexible since his translation to categories
requires the desired category to only implement primitives
actually used for translation of a particular term at
hand.

Conor McBride's ``Crude but effective stratification''
\cite{McBride:Crude2011, McBride:Crude2012} (systemic
treatment of this idea is given by Favonia and others
\cite{Favonia:Displacement2023}) introduces an
eye-opening idea of using \textit{displacement} as a
lightweight form of universe polymorphism: to reuse term
typed at a lower universe level, lift all levels in a
term by adding a constant shift $s$. From the point of
view advocated by this paper, you can see it as an act
of reinterpreting an expression in an internal model of
a type system where the universe at level $i$ is interpreted
as a universe at level $i+s$
(Section~\ref{twoLTT-displacement}).

In ``Parametricity and Dependent Types'' by
Jean-Philippe Bernardy, Patrik Jansson and Ross Paterson
\cite{Bernardy:Parametricity2010}, authors introduce a
direct translation procedure to yield parametricity
theorems and proofs from polymorphic definitions in a
dependently typed system. The definition of <<external
parametricity>> we follow, though, is given in
``Internal Parametricity, without an Interval'' by
Thorsten Altenkirch, Yorgo Chamoun, Ambrus Kaposi and
Michael Shulman \cite{Altenkirch:Parametricity2024}.
Once again, we can reformulate the translation they give
as reinterpretation in an internal model. Altenkirch and
others noted an issue with this approach: if we apply
the external translation to the judgment with nonempty
context, context becomes parametric, too; that's the
reason they add \texttt{rel} operation and yield type
system with \textit{computational properties of higher
observational type theory}. We sidestep the issue with
context since reinterpretation is applied only to terms
whose context is also <<reinterpretable>> and is not
affected by the translation
(Section~\ref{twoLTT-parametricity}).

\subsection{What we could've done differently}

\subsubsection{Just use metaprogramming instead}

In a recent paper ``A Certifying Proof Assistant for
Synthetic Mathematics in Lean'', Wojciech Nawrocki and
others \cite{Nawrocki:Lean2026} adapt \textit{Compiling
to categories} approach \cite{Elliott:Compcat2017} to
Lean's system with dependent types, using
Categories-with-Families to represent internal models
and employing Lean's powerful metaprogramming facilities
instead of a compiler plugin. In doing so, they are able
to provide a framework for doing synthetic mathematics
in Lean, extracting proofs done in MLTT applied to more
specific categories. Thus our results in
Section~\ref{twoLTT-synthetic} are close to the results shown in
their paper. To add to that, their work is already
implemented in a practical proof assistant. Yet our type
system involves no metaprogramming since
reinterpretation is internal to it, and it handles
conversion rule. In addition, the focus is a bit
different, where we use reinterpretation to express
other polymorphic features of a language.

\subsubsection{
    Computational modality of global elements}
\label{global-elements}

Since we introduced a family of modalities of
reinterpretable terms $\Box_\ell$ in
Section~\ref{twoLTT}, one could hope to allow computation in
boxed terms. Generalizing a bit, the following
definition will be useful further:

\begin{definition}
    A dependently-typed system with a type of functions
    $(\to)$ has a \textit{computational modality of
    global elements} $\Box$ if there is a one-hole term
    $\mathrm{nec}\,\_$ and an atomic two-hole term
    $\mathrm{K}\,\_\,\_$ such that the following rules
    are admissible:
    \begin{mathpar}
    \inferrule{\cdot\vdash T\,\mathrm{type}}
        {\Gamma\vdash\Box\,T\,\mathrm{type}}
    \and
    \inferrule{\cdot\vdash t:T}
        {\Gamma\vdash\mathrm{nec}\,t:\Box\,T}
    \and
    \inferrule{
        \Gamma\vdash f:\Box\,(S\to T)\\
        \Gamma\vdash t:\Box\,S
    }{\Gamma\vdash K\,f\,t:\Box\,T}
    \end{mathpar}
    And, in addition, the following reduction is
    admissible: $K\,(\mathrm{nec}\,\lambda x.t)\,
    (\mathrm{nec}\,s)\leadsto^*
    \mathrm{nec}(t[s/x])$.
\end{definition}

With an exception of a computational rule, these are
exactly the properties of modalities $\Box_\ell$ from
the original Section~\ref{twoLTT}. The intuition behind
the reduction is that, if we apply a ``global function''
to a ``global argument'', we then should be able to
``inline'' this application, what this reduction
precisely does.

Then, a corresponding weakening of computational
reinterpretation can be defined as follows.

\begin{definition}[Modal computational reinterpretation]
    We call a type operator $F[\_]$ a modal
    computational reinterpretation of types if it is
    defined on boxed types (if
    $\Gamma\vdash\Box\,T\,\mathrm{Type}$, then
    $\Gamma\vdash F[T]\,\mathrm{Type}$), respects
    computation on types (if $T\to T'$, then
    $F[T]\to^* F[T']$) and computes as a primitive
    recursive function on normal forms of types.
    Analogously, a term with a hole
    $\llbracket\_\rrbracket$ is a modal computational
    reinterpretation of terms if it respects computation
    on terms, computes as a primitive recursive
    function on normal forms of terms and, if
    $\Gamma\vdash t:\Box\,T$, then
    $\Gamma\vdash\llbracket t\rrbracket:F[T]$.
\end{definition}

Maybe this definition of reinterpretation is ``nice''?
Regrettably, there's a no-go theorem which is similar to
na\"{\i}ve version \ref{no-go} which holds in this case:

\begin{theorem}[No-go theorem of modal reinterpretation]
    In a confluent type system with modal computational
    reinterpretation, the following should hold for
    arbitrary well-typed $s$ and $t$:
    \[
        \mathtt{app}\llbracket\mathrm{nec}\,\lambda x.t
        \rrbracket\llbracket\mathrm{nec}\,s\rrbracket
        \equiv\llbracket\mathrm{nec}(t[s/x])
        \rrbracket
    \]
    where $\equiv$ is definitional equality and
    $\mathtt{app}$ is how $\llbracket\_\rrbracket$
    acts on $K$.
\end{theorem}

\begin{proof}
    By direct example. We are given well-formed types
    $\cdot\vdash S,T\,\mathrm{type}$ and well-typed
    terms $\cdot\vdash s:S$ and $\cdot,x:S\vdash t:T$.
    Consider the following term:
    \[
        \Gamma,x:\Box(S\to T),y:\Box\,S\vdash
            \llbracket K\,x\,y\rrbracket:F[T]
    \]
    Since $K$ is atomic, the term under interpretation is in
    normal form; hence the whole term computes to
    $\mathtt{app}\llbracket x\rrbracket \llbracket y
    \rrbracket$. Substituting $\mathrm{nec}\,
    \lambda x.t$ for $x$ and $\mathrm{nec}\,s$ for $y$,
    we get $\mathtt{app}\llbracket\mathrm{nec}\,
    \lambda x.t\rrbracket\llbracket\mathrm{nec}\,
    s\rrbracket$.

    On the other hand, if we perform substitution first,
    we would get $\Gamma\vdash\llbracket K\,
    (\mathrm{nec}\,\lambda x.t)(\mathrm{nec}\,s)
    \rrbracket:F[T]$. Since modal reinterpretation
    respects reduction, this reduces to
    $\llbracket\mathrm{nec}(t[s/x])\rrbracket$.

    In a confluent type system with functions,
    multi-step reduction commutes with substitution;
    therefore, two terms are definitionally equivalent.
\end{proof}

This, of course, doesn't prohibit us from using such a
modality: as is with na\"{\i}ve interpretation, this
approach works if we manage to constrain models to those
where equalities hold definitionally.

Another benefit of modal computational reinterpretation
is the following property:

\begin{theorem}[Collapse of computational modalities]
    The following implication about $\twoLTT$ is true:
    \[\dfrac{
        \exists\kappa:
        \Gamma.\,{}_\ell\text{\faLock}\vdash_\kappa
        t:T
    }{\forall j\geqslant\mathrm{level}(T):
        \Gamma.\,{}_\ell\text{\faLock}\vdash_j
    \mathrm{nf}(t):\mathrm{nf}(T)}\]
\end{theorem}

We do not give a direct proof here, but the general idea
is that any level higher than the level of $T$ can only
appear in some redex of $t$ or $T$. Then, if modalities
were computational, it would be enough to use a single
modality $\Box$ instead of using a family of modalities
$\Box_\ell$: level of type under box would tell how many
universe models we need to provide.

But, again, it all comes down to defining the type of
\textit{definitional} internal models.

\subsubsection{Interpreting interpretation}

In Section~\ref{twoLTT}, we introduced our
``reinterpretable terms'' modality in a very
constraining way specifically so that we wouldn't have
to provide interpretation of this modality inside
internal models. Once again, the problem is with terms
that try to interpret variables from CwF context;
the simplest example is $\llbracket \lambda x.\llbracket x
\rrbracket_1\rrbracket_2$. In order to be able to
interpret such a term, we would have to define \textit{
what reinterpretation operator means inside an internal
model}. This quickly gets out of hand: rather than just
turning the type of models into a (recursive? inductive?)
type, the type of models would become such a type that a
type of one of its fields depends on the value of the
whole model. Actually, this happened before:
Kolomatskaia et al. \cite{Kolomatskaia_Shulman_2025}
have already defined a type like this inside their
\textit{displayed type theory}: type SST of
semi-simplicial types is defined coinductively in terms
of values of itself\footnote{More generally, Shulman
dubs such types \textit{modal codatatypes}
\cite{shulman2026}.}. While defining type of models in
a similar vein could be an interesting development of the type
system presented in this paper, we deem it out of scope
for now.

\section{Conclusions and further work}

In retrospect, we managed to define two
self-interpreting type systems, $\SSTLC$ and $\twoLTT$,
meaning that they have a notion of their own ``internal
models'' which they then can use to reinterpret their
own terms inside such models. In doing so, such type
systems get new features which previously were defined
as separate builtin quirks of a system. Basic
metatheoretical properties are proven; for future work
we have the following plans:

\begin{itemize}
    \item We might try making locking discipline in
        Section~\ref{SSTLC} more flexible, allowing
        more than one lock in a context. Most likely,
        this would require us to incorporate modalities
        in the definition of an ``internal model'' type
        and bring us closer to the truly reflective
        system.
    \item While our technique to reinterpret conversion
        rule in Section~\ref{twoLTT} works, a current
        approach to equality witness extraction is a
        \textit{hack}. A proper solution would require us
        to specify witness extraction as a redex and,
        possibly, incorporate certain \textit{algebras
        of reduction strategies} to make this process
        more controllable in proof-relevant settings.
    \item A Categories-with-Families approach to
        specifying internal models here is, in a sense,
        a ``deep'' interface: while all of its
        facilities are properly used in interpretation,
        we have a low-hanging fruit of being able to
        interpret terms in arbitrary CwF context. This
        hints at a direct extension of our modality of
        reinterpretable terms into full-blown contextual
        types \cite{Nanevski:2008:CMT}.
    \item Having implemented contextual types, it then
        becomes possible to handle ``partial''
        reinterpretation of terms: if a term under
        interpretation doesn't use a certain feature of
        a type system, then you shouldn't have to
        specify it in a model! This would bring our
        language closer to Conal Elliott's approach of
        ``Compiling to Categories''
        \cite{Elliott:Compcat2017} and closer to
        an \textit{actual} type system for synthetic
        mathematics.
    \item It would be nice to make equality in $\twoLTT$
        more extensional, most probably by making the type
        system closer to Homotopy type theory
        \cite{hottbook}.
    \item At some moment we would have to pay due
        diligence: provide mechanized proofs of
        metatheoretical properties of systems and an
        implementation of $\twoLTT$ in a practical proof
        assistant.
\end{itemize}

\printbibliography
\appendix
\section{A type of internal models of
\texorpdfstring{$\twoLTT$}{[[2]]LTT}}
\label{twoltt-internal-model-def}

A comprehension $[x:A,y:B,z:C]$ denotes
a right-nested $\Sigma$-type. $\{x:A\}\to B$ is a
syntactic sugar meaning that argument $x$ is always
inferred from $B$ rather than provided explicitly.
$\{x\}\to B$ means that type of $x$ is inferred from
$B$. Transports along previously introduced equalities
are omitted. Since our language doesn't yet have
$\eta$-equivalence, uniqueness properties are skipped.

\[\mathcal{M}_{\ell_0,\ldots,\ell_\kappa}^{
    \ell'_0,\ldots,\ell'_\kappa}={[}\]

\subsection{Base CwF part}

\[\begin{array}{l}
    \mathcal{C}:\mathcal{U}_{\max(\ell'_0,\ldots,
        \ell'_\kappa)},\\
    \triangleright:\mathcal{C}\to\mathcal{C}\to
        \mathcal{U}_{\max(\ell_0,\ldots,\ell_\kappa)},\\
    \mathrm{id}:\{\Gamma:\mathcal{C}\}\to
        \Gamma\triangleright\Gamma,\\
    \circ:\{\Gamma\,\Delta\,\Xi\}\to
        (\Delta\triangleright\Xi)\to
        (\Gamma\triangleright\Delta)\to
        (\Gamma\triangleright\Xi),\\
    \mathrm{id-}\circ:\{\Gamma,\Delta\}\to\{
        \sigma:\Gamma\to\Delta\}\to
        \mathrm{id}\circ\sigma=\sigma,\\
    \circ\mathrm{-id}:\{\Gamma,\Delta\}\to\{
        \sigma:\Gamma\to\Delta\}\to
        \sigma\circ\mathrm{id}=\sigma,\\
    \circ\mathrm{-assoc}:\{\rho\,\sigma\,\tau\}\to
        \rho\circ(\sigma\circ\tau)=
        (\rho\circ\sigma)\circ\tau,\\
    \cdot:\mathcal{C},\\
    !:\{\Gamma\}\to\Gamma\triangleright\cdot,\\
\end{array}\]

\subsection{Base type formers for each
\texorpdfstring{$j$}{j}}

\[\begin{array}{l}
    \mathrm{Ty}_j:\mathcal{C}\to\mathcal{U}_{\ell'_j},
    \\
    \_[\_]:\{\Gamma\,\Delta\}\to\mathrm{Ty}_j\Gamma\to
        \Delta\triangleright\Gamma\to
            \mathrm{Ty}_j\Delta,\\
    \mathrm{sub/id}:\{T\}\to T[\mathrm{id}]=T,\\
    \mathrm{sub/comp}:\{T\,\sigma\,\tau\}\to
        T[\sigma\circ\tau]=T[\sigma][\tau],\\
    \mathrm{Tm}_j:(\Gamma:\mathcal{C})\to
        \mathrm{Ty}_j\,\Gamma\to\mathcal{U}_{\ell_j},\\
    \_[\_]:\{\Gamma\,\Delta\,T\}\to
        \mathrm{Tm}_j\,\Gamma\,T\to(\sigma:
        \Delta\triangleright\Gamma)\to
        \mathrm{Tm}_j\,\Delta\,T[\sigma],\\
    \mathrm{sub/id}:\{T\}\to\{t:\mathrm{Tm}\_\,T\}\to
        t[\mathrm{id}]=t,\\
    \mathrm{sub/comp}:\{T\,\sigma\,\tau\}\to
        \{t:\mathrm{Tm}_j\,\_\,T\}\to
        t[\sigma\circ\tau]=t[\sigma][\tau],\\
    (.):(\Gamma:\mathcal{C})\to\mathrm{Ty}_j\,\Gamma\to
        \mathcal{C},\\
    p:\{\Gamma\,T\}\to\Gamma.T\triangleright\Gamma,\\
    q:\mathrm{Tm}\,\Gamma.T\,T[p],\\
    \langle,\rangle:\{\Delta\,\Gamma\}\to
        (\sigma:\Delta\triangleright\Gamma)\to
        \mathrm{Tm}\,\Delta\,T[\sigma]\to
        \Delta\triangleright\Gamma.T,\\
    \mathrm{ext/l}:\{\sigma\,t\}\to
        p\circ\langle\sigma, t\rangle=\sigma,\\
    \mathrm{ext/r}:\{\sigma\,t\}\to
        q[\langle \gamma, t\rangle]=t,\\
    \mathrm{Nat}_j:\{\Gamma\}\to\mathrm{Ty}_j\,
        \Gamma,\\
    0:\{\Gamma\}\to\mathrm{Tm}_j\,\Gamma\,
        \mathrm{Nat}_j,\\
    S:\{\Gamma\}\to\mathrm{Tm}_j\,\Gamma\,
        \mathrm{Nat}_j\to\mathrm{Tm}_j\,\Gamma
            \,\mathrm{Nat}_j,\\
    \mathrm{Id}:\{\Gamma\}\to(T:\mathrm{Ty}_j\,\Gamma)
        \to\mathrm{Tm}_j\,\Gamma\,T
        \to\mathrm{Tm}_j\,\Gamma\,T
        \to\mathrm{Ty},\\
    \mathrm{refl}:\{\Gamma\,T\}\to
        \{t\,t':\mathrm{Tm}_j\,\Gamma\,T\}\to
        t=t'\to\mathrm{Tm}_j\,\Gamma\,
            (\mathrm{Id}\,T\,t\,t'),\\
    \mathrm{Nat}[\,] : \{\sigma\}\to
        \mathrm{Nat}_j[\sigma]=\mathrm{Nat}_j,\\
    0[\,] : 0[\sigma]=0,\\
    \mathrm S[\,] :
        (\mathrm S\,t)[\sigma]=\mathrm S\,(t[\sigma]),\\
    \mathrm{Id}[\,] :
        (\mathrm{Id}\,T\,t\,t')[\sigma] =
        \mathrm{Id}\,(T[\sigma])\,(t[\sigma])\,
        (t'[\sigma]),\\
    \mathrm{refl}[\,] :
        (\mathrm{refl}\,e)[\sigma] =
        \mathrm{refl}\,(\mathrm{ap}_{\_[\sigma]}\,e),
\end{array}\]

\subsection{Universes, for each
\texorpdfstring{$j<\kappa$}{j<k}}
\[\begin{array}{l}
    \mathrm U_j :
        \{\Gamma\}\to\mathrm{Ty}_{j+1}\,\Gamma,\\
    \mathrm U[\,] : \mathrm U_j[\sigma]=\mathrm U_j,\\
    \mathrm{El} :
        \{\Gamma\}\to\mathrm{Tm}_{j+1}\,
        \Gamma\,\mathrm U_j\to\mathrm{Ty}_j\,\Gamma,\\
    \mathrm{El}[\,] :
        (\mathrm{El}\,t)[\sigma]=
        \mathrm{El}\,(t[\sigma]),\\
    \mathrm u_j : \{\Gamma\}\to\mathrm{Tm}_{j+2}\,
        \Gamma\,\mathrm U_{j+1}\ \ (j+1<\kappa),\\
    \mathrm{El/u} : \mathrm{El}\,\mathrm u_j =
        \mathrm U_j,\\
    \mathrm{nat}_j : \{\Gamma\}\to\mathrm{Tm}_{j+1}\,
        \Gamma\,\mathrm U_j,\\
    \mathrm{El/nat} :
        \mathrm{El}\,\mathrm{nat}_j=\mathrm{Nat}_j,\\
    \mathrm{id}_j :
        \{\Gamma\}\to(a:\mathrm{Tm}\,\Gamma\,
        \mathrm U_j)\to\mathrm{Tm}\,\Gamma\,(
        \mathrm{El}\,a)\to\mathrm{Tm}\,\Gamma\,(
        \mathrm{El}\,a)\to\mathrm{Tm}\,\Gamma\,
        \mathrm U_j,\\
    \mathrm{El/id} :
        \mathrm{El}\,(\mathrm{id}_j\,a\,t\,t')=
        \mathrm{Id}\,(\mathrm{El}\,a)\,t\,t',
\end{array}\]
together with naturality of $\mathrm u_j,\mathrm{nat}_j,\mathrm{id}_j$ under substitution.

\subsection{\texorpdfstring{$\Pi$}{Pi},
\texorpdfstring{$\Sigma$}{Sigma} and eliminators, for
each \texorpdfstring{$i,j\le\kappa$}{i,j<=k}}
\[\begin{array}{l}
    \Pi : \{\Gamma\}\to(A:\mathrm{Ty}_i\,\Gamma)
    \to\mathrm{Ty}_j\,(\Gamma.A)
    \to\mathrm{Ty}_{i\sqcup j}\,\Gamma,\\
    \Pi[\,] : (\Pi\,A\,B)[\sigma]=
    \Pi\,(A[\sigma])\,(B[\sigma^+]),\\
    \mathrm{lam} :
    \mathrm{Tm}_j\,(\Gamma.A)\,B\to
    \mathrm{Tm}_{i\sqcup j}\,\Gamma\,(\Pi\,A\,B),\\
    \mathrm{app} :
    \mathrm{Tm}\,\Gamma\,(\Pi\,A\,B)
    \to(a:\mathrm{Tm}\,\Gamma\,A)
    \to\mathrm{Tm}_j\,\Gamma\,(B[\overline a]),\\
    \Pi\text{-}\beta : \mathrm{app}\,(\mathrm{lam}\,b)\,
        a = b[\overline a],\\
    \mathrm{lam}[\,] :
        (\mathrm{lam}\,b)[\sigma]=
        \mathrm{lam}\,(b[\sigma^+]),\\
    \mathrm{app}[\,] :
        (\mathrm{app}\,f\,a)[\sigma]=
        \mathrm{app}\,(f[\sigma])\,(a[\sigma]),\\
    \mathrm{pi} : (a:\mathrm{Tm}\,\Gamma\,\mathrm U_i)
    \to\mathrm{Tm}\,(\Gamma.\mathrm{El}\,a)\,
    \mathrm U_j\to\mathrm{Tm}\,\Gamma\,
    \mathrm U_{i\sqcup j}\ (i,j<\kappa),\\
    \mathrm{El/pi} : \mathrm{El}\,(\mathrm{pi}\,a\,b)=
    \Pi\,(\mathrm{El}\,a)\,(\mathrm{El}\,b),\\
    \Sigma : \{\Gamma\}\to(A:\mathrm{Ty}_i\,\Gamma)
    \to\mathrm{Ty}_j\,(\Gamma.A)
    \to\mathrm{Ty}_{i\sqcup j}\,\Gamma,\\
    \Sigma[\,] : (\Sigma\,A\,B)[\sigma]=
    \Sigma\,(A[\sigma])\,(B[\sigma^+]),\\
    \mathrm{pair} : (a:\mathrm{Tm}\,\Gamma\,A)\to
        \mathrm{Tm}\,\Gamma\,(B[\overline a])\to
        \mathrm{Tm}\,\Gamma\,(\Sigma\,A\,B),\\
    \mathrm{fst} : \mathrm{Tm}\,\Gamma\,(\Sigma\,A\,B)
        \to\mathrm{Tm}\,\Gamma\,A,\\
    \mathrm{snd} : (s:\mathrm{Tm}\,\Gamma\,
        (\Sigma\,A\,B))\to\mathrm{Tm}\,\Gamma\,
        (B[\overline{\mathrm{fst}\,s}]),\\
    \Sigma\text{-}\beta_1 :
        \mathrm{fst}\,(\mathrm{pair}\,a\,b)=a,\\
    \Sigma\text{-}\beta_2 :
        \mathrm{snd}\,(\mathrm{pair}\,a\,b)=b,\\
    \mathrm{sigma},\ \mathrm{El/sigma},
    \ \mathrm{pair}[\,],\ \mathrm{fst}[\,],
    \ \mathrm{snd}[\,]\ \text{analogous},\\
    \mathrm{natind} :
        (M:\mathrm{Ty}_i\,(\Gamma.\mathrm{Nat}_j))
        \to\mathrm{Tm}\,\Gamma\,(M[\overline 0])
        \to\mathrm{Tm}\,(\Gamma.\mathrm{Nat}_j.M)\,
        (M[\langle p\circ p,\ \mathrm S\,(q[p])\rangle])
        \to\\\qquad
        \to(n:\mathrm{Tm}\,\Gamma\,\mathrm{Nat}_j)
        \to\mathrm{Tm}_i\,\Gamma\,(M[\overline n]),\\
    \mathrm{Nat}\text{-}\beta_0 :
        \mathrm{natind}\,M\,z\,s\,0=z,\\
    \mathrm{Nat}\text{-}\beta_{\mathrm S} :
        \mathrm{natind}\,M\,z\,s\,(\mathrm S\,n)=
        s[\langle\overline n,\ \mathrm{natind}\,M\,z\,s
        \,n\rangle],\\
    \mathrm{natind}[\,],\\
    \mathrm J : (T:\mathrm{Ty}_j\,\Gamma)
        (t:\mathrm{Tm}\,\Gamma\,T)
        (M:\mathrm{Ty}_i\,(\Gamma.T.
            \mathrm{Id}\,(T[p])\,(t[p])\,q))
        \to\mathrm{Tm}\,\Gamma\,
            (M[\langle\overline t,\mathrm{refl}\rangle])
            \to\\
        \qquad\to(t':\mathrm{Tm}\,\Gamma\,T)
        (e:\mathrm{Tm}\,\Gamma\,(\mathrm{Id}\,T\,t\,t'))
        \to\mathrm{Tm}_i\,\Gamma\,
        (M[\langle\overline{t'},e\rangle]),\\
    \mathrm{Id}\text{-}\beta :
        \mathrm J\,T\,t\,M\,r\,t\,\mathrm{refl}=
        r,\\ \mathrm J[\,],
\end{array}\]

\subsection{End of definition of a model type}

\[{]}\]

\section{Interpretation procedures}
\label{interpretation-definition}

\paragraph{Contexts and variables.}
\begin{align*}
m\llbracket\cdot\rrbracket &= m.\cdot &
m\llbracket\Delta, x:T\rrbracket &=
    m\llbracket\Delta\rrbracket\,.\,
    m\llbracket\Delta\vdash T\rrbracket\\
m\llbracket\Delta,x:T\vdash x\Rightarrow\rrbracket&=
    (q,\; T) &
m\llbracket\Delta, y:S \vdash x\Rightarrow\rrbracket &=
    (u[p],\; T)\ \text{ where }
(u,T)=m\llbracket\Delta\vdash x\Rightarrow\rrbracket
\end{align*}
(the second component is the syntactic $T$; its
interpretation is weakened by $\mathrm{wk}_m$, see
below).

\paragraph{Types.}
\begin{align*}
m\llbracket\Delta\vdash \mathcal U_j\rrbracket &=
    \mathrm U_j \\
m\llbracket\Delta\vdash \mathrm{Nat}_j\rrbracket &=
    \mathrm{Nat}_j\\
m\llbracket\Delta\vdash \mathrm{El}(t)\rrbracket &=
    \mathrm{El}\big(\pi_1\,m\llbracket\Delta\vdash
        t\Rightarrow\rrbracket\big) \\
m\llbracket\Delta\vdash t_1 =_T t_2\rrbracket &=
\mathrm{Id}\big(m\llbracket\Delta\vdash T\rrbracket,\,
m\llbracket\Delta\vdash t_1\Leftarrow T\rrbracket,\,
m\llbracket\Delta\vdash t_2\Leftarrow T\rrbracket\big)\\
m\llbracket\Delta\vdash (x:S)\times T\rrbracket &=
\Sigma\big(m\llbracket\Delta\vdash S\rrbracket,
\ m\llbracket\Delta,x:S\vdash T\rrbracket\big)\\
m\llbracket\Delta\vdash (x:S)\to T\rrbracket &=
\Pi\big(m\llbracket\Delta\vdash S\rrbracket,
\ m\llbracket\Delta,x:S\vdash T\rrbracket\big)
\end{align*}
$\square_\kappa T$ and $t\llbracket T\rrbracket$ cannot occur under a lock and need no clause.

\paragraph{Checked terms.} Whenever a rule demands a type of a particular shape, the type is first brought to weak head normal form by the fixed strategy and a transport along $\mathrm{eq}_m$ is inserted; we write $t \Leftarrow^{\!*} T$ for this.
\begin{align*}
m\llbracket\Delta\vdash 0 \Leftarrow^{\!*} \mathrm{Nat}_j\rrbracket &= 0 \\
m\llbracket\Delta\vdash \mathrm S\,t \Leftarrow^{\!*} \mathrm{Nat}_j\rrbracket &= \mathrm S\,m\llbracket\Delta\vdash t\Leftarrow \mathrm{Nat}_j\rrbracket\\
m\llbracket\Delta\vdash (s,t) \Leftarrow^{\!*} (x:S)\times T\rrbracket &= \mathrm{pair}\big(m\llbracket\Delta\vdash s\Leftarrow S\rrbracket,\ m\llbracket\Delta\vdash t\Leftarrow T[(s:S)/x]\rrbracket\big)\\
m\llbracket\Delta\vdash \lambda x.t \Leftarrow^{\!*} (x:S)\to T\rrbracket &= \mathrm{lam}\big(m\llbracket\Delta,x:S\vdash t\Leftarrow T\rrbracket\big)\\
m\llbracket\Delta\vdash \mathrm{refl} \Leftarrow^{\!*} t_1 =_T t_2\rrbracket &= \mathrm{refl}\big(\mathrm{eq}_m(\Delta; t_1, t_2)\big)\\
m\llbracket\Delta\vdash t \Leftarrow T\rrbracket &= \mathrm{tr}_{\mathrm{Tm}\,m\llbracket\Delta\rrbracket}\big(\mathrm{eq}_m(\Delta; S, T),\ u\big) \\
   &\text{where } (u,S)=m\llbracket\Delta\vdash t\Rightarrow\rrbracket \quad(\textsc{Check-Inferred})
\end{align*}

\paragraph{Synthesising terms.} Below $(u,S)=m\llbracket\Delta\vdash t\Rightarrow\rrbracket$, $\mathrm{whnf}(S)$ is matched against the required shape, and $u^{\dagger}$ denotes $u$ transported along $\mathrm{eq}_m(\Delta;S,\mathrm{whnf}(S))$.
\begin{align*}
m\llbracket\Delta\vdash \mathtt U_\kappa\Rightarrow\rrbracket &= (\mathrm u_\kappa,\ \mathcal U_{\kappa+1}) \\
m\llbracket\Delta\vdash \mathtt{Nat}_\kappa\Rightarrow\rrbracket &= (\mathrm{nat}_\kappa,\ \mathcal U_\kappa)\\
m\llbracket\Delta\vdash \mathtt{Pi}(x:s)t\Rightarrow\rrbracket &= \big(\mathrm{pi}(a, b),\ \mathcal U_{j\sqcup\kappa}\big)\\
 &\text{with } (a,\mathcal U_j) \leftarrow m\llbracket\Delta\vdash s\Rightarrow\rrbracket,\ (b,\mathcal U_\kappa) \leftarrow m\llbracket\Delta,x:\mathrm{El}(s)\vdash t\Rightarrow\rrbracket\\
m\llbracket\Delta\vdash \mathtt{Sigma}(x:s)t\Rightarrow\rrbracket &= \big(\mathrm{sigma}(a,b),\ \mathcal U_{j\sqcup\kappa}\big) \\ &\text{likewise}\\
m\llbracket\Delta\vdash \mathtt{Id}\,t\,t_1\,t_2\Rightarrow\rrbracket &= \big(\mathrm{id}(a, e_1, e_2),\ \mathcal U_\kappa\big)\\
&(a,\mathcal U_\kappa) \leftarrow m\llbracket\Delta\vdash t\Rightarrow\rrbracket,\ e_i = m\llbracket\Delta\vdash t_i\Leftarrow\mathrm{El}(t)\rrbracket\\
m\llbracket\Delta\vdash t\,s\Rightarrow\rrbracket &= \big(\mathrm{app}(u^\dagger, a)^{\,\mathrm{sub}},\ T[(s:S')/x]\big)
 \\ &\mathrm{whnf}(S)=(x:S')\to T,\ a=m\llbracket\Delta\vdash s\Leftarrow S'\rrbracket\\
m\llbracket\Delta\vdash \pi_1 t\Rightarrow\rrbracket &= (\mathrm{fst}\,u^\dagger,\ S') \\ &\mathrm{whnf}(S)=(x:S')\times T\\
m\llbracket\Delta\vdash \pi_2 t\Rightarrow\rrbracket &= \big((\mathrm{snd}\,u^\dagger)^{\mathrm{sub}},\ T[(\pi_1 t)/x]\big)\\
m\llbracket\Delta\vdash \mathrm{natind}_{x.T}\,t\,\{s_0,s_{\mathrm S}\}\Rightarrow\rrbracket &=
 \big(\mathrm{natind}(M, z, s, u^\dagger)^{\mathrm{sub}},\ T[t/x]\big)\\
 & M = m\llbracket\Delta,x:\mathrm{Nat}_\kappa\vdash T\rrbracket,\ z = m\llbracket\Delta\vdash s_0\Leftarrow T[0/x]\rrbracket,\\
 & s = m\llbracket\Delta,y:\mathrm{Nat}_\kappa, z:T[y/x]\vdash s_{\mathrm S}\,y\,z \Leftarrow T[\mathrm S y/x]\rrbracket\\
m\llbracket\Delta\vdash \mathrm J_{xy.S}\,t\,s\Rightarrow\rrbracket &= \big(\mathrm J(A, e_1, M, r, e_2, u^\dagger)^{\mathrm{sub}},\ S[(t_2:T)/x, t/y]\big)
 \\ & \mathrm{whnf}(S_t) = t_1 =_T t_2,\ A = m\llbracket\Delta\vdash T\rrbracket,\ e_i = m\llbracket\Delta\vdash t_i\Leftarrow T\rrbracket,\\
 & M = m\llbracket\Delta,x:T,y:t_1=_Tx\vdash S\rrbracket,\ r = m\llbracket\Delta\vdash s\Leftarrow S[t_1/x,\mathrm{refl}/y]\rrbracket\\
m\llbracket\Delta\vdash (t:T)\Rightarrow\rrbracket &= \big(m\llbracket\Delta\vdash t\Leftarrow T\rrbracket,\ T\big)
\end{align*}
The superscript $(\_)^{\mathrm{sub}}$ is a transport along the \emph{substitution lemma}
\[
\mathrm{sub}_m :\ m\llbracket\Delta,x:S\vdash T\rrbracket[\overline{m\llbracket\Delta\vdash s\Leftarrow S\rrbracket}] \;=\; m\llbracket\Delta\vdash T[(s:S)/x]\rrbracket
\]
(and its weakening instance $\mathrm{wk}_m$), proved by induction on $T$ from the naturality equations $\mathrm{sub/id}$, $\mathrm{sub/comp}$, $\Pi[\,]$, $\Sigma[\,]$, $\mathrm{Id}[\,]$, $\mathrm{El}[\,]$, $\mathrm{U}[\,]$, $\mathrm{Nat}[\,]$ of the model; it is a closed meta-level term and is computed, not postulated.

\paragraph{Equality witnesses.} For $a \equiv b$ (both types, or both terms of the same type), with $\leadsto_{\mathrm s}$ the deterministic strict strategy:
\[
\mathrm{eq}_m(\Delta; a, b) =
\begin{cases}
\mathrm{refl} & a =_\alpha b,\\
\mathrm{trans}\big(\mathrm{step}_m(\Delta; a\leadsto_{\mathrm s} a'),\ \mathrm{eq}_m(\Delta; a', b)\big) & a \leadsto_{\mathrm s} a',\\
\mathrm{sym}\,\mathrm{eq}_m(\Delta; b, a) & a \text{ normal}, b \text{ not}.
\end{cases}
\]
$\mathrm{step}_m$ maps the contracted redex to the corresponding model equation — $\Pi$-$\beta$, $\Sigma$-$\beta_i$, $\mathrm{Nat}$-$\beta_{0/\mathrm S}$, $\mathrm{Id}$-$\beta$, the $\mathrm{El}$-equations, $\mathrm{refl}$ for $(t:T)\leadsto t$ — composed with $\mathrm{sub}_m$ where the redex substitutes, and placed under the surrounding term formers by $\mathrm{ap}$ (interpretation of a former is a meta-level function of the interpretations of its immediate subterms, so congruence is $\mathrm{ap}$). Termination: the argument of the recursive call is strictly closer to normal form, and $a,b$ are MLTT terms, hence strongly normalising; well-definedness: convertible normal MLTT terms are $\alpha$-equal (confluence; no $\eta$).

Note a nice property: under the identity model
$I_\ell$ all equations used by $\mathrm{step}_m$ and
$\mathrm{sub}_m$ are $\mathrm{refl}$, so all transports
vanish and $I_\ell\llbracket\Delta\vdash t\Leftarrow T
\rrbracket$ is the standard shallow embedding of $t$.

\section{Definition of a displacement model}
\label{twoltt-displacement-model-def}

\[\begin{array}{ll}
    \uparrow^{\kappa}_{\ell} = &[\;
\mathcal C = \mathcal U_{\ell+\kappa},\quad
\Gamma\triangleright\Delta = \Gamma\to\Delta,\quad
\mathrm{id}=\lambda x.x,\quad f\circ g=\lambda x.f(gx),\quad
\cdot = \top,\quad {!}=\lambda\_.(),\\
&\mathrm{Ty}_j\,\Gamma=\Gamma\to\mathcal U_{j+\kappa},\quad T[\sigma]=\lambda\gamma.T(\sigma\gamma),\quad
\mathrm{Tm}_j\,\Gamma\,T=(\gamma:\Gamma)\to T\gamma,\quad t[\sigma]=\lambda\gamma.t(\sigma\gamma),\\
&\Gamma.T=(\gamma:\Gamma)\times T\gamma,\quad p=\pi_1,\quad q=\pi_2,\quad \langle\sigma,t\rangle=\lambda\delta.(\sigma\delta,t\delta),\\
&\mathrm{Nat}_j=\lambda\_.\mathrm{Nat}_{j+\kappa},\quad 0=\lambda\_.0,\quad \mathrm S\,t=\lambda\gamma.\mathrm S(t\gamma),\\
&\mathrm{natind}\,M\,z\,s\,n=\lambda\gamma.\mathrm{natind}_{x.M(\gamma,x)}(n\gamma)\{z\gamma,\ \lambda y\,r.\,s((\gamma,y),r)\},\\
&\mathrm{Id}\,T\,t\,t'=\lambda\gamma.(t\gamma=_{T\gamma}t'\gamma),\quad \mathrm{refl}\,e=\lambda\gamma.\mathrm{tr}(e,\mathrm{refl}),\quad
\mathrm J\,T\,t\,M\,r\,t'\,e=\lambda\gamma.\mathrm J_{xy.M((\gamma,x),y)}(e\gamma)(r\gamma),\\
&\mathrm U_j=\lambda\_.\mathcal U_{j+\kappa},\quad \mathrm{El}\,t=\lambda\gamma.\mathrm{El}(t\gamma),\quad
\mathrm u_j=\lambda\_.\mathtt U_{j+\kappa},\quad \mathrm{nat}_j=\lambda\_.\mathtt{Nat}_{j+\kappa},\\
&\mathrm{pi}\,a\,b=\lambda\gamma.\mathtt{Pi}(x:a\gamma)\,b(\gamma,x),\quad
\mathrm{sigma}\,a\,b=\lambda\gamma.\mathtt{Sigma}(x:a\gamma)\,b(\gamma,x),\\
&\mathrm{id}_j\,a\,t\,t'=\lambda\gamma.\mathtt{Id}\,(a\gamma)\,(t\gamma)\,(t'\gamma),\\
&\Pi\,A\,B=\lambda\gamma.(x:A\gamma)\to B(\gamma,x),\quad \mathrm{lam}\,b=\lambda\gamma\,x.b(\gamma,x),\quad \mathrm{app}\,f\,a=\lambda\gamma.f\gamma(a\gamma),\\
&\Sigma\,A\,B=\lambda\gamma.(x:A\gamma)\times B(\gamma,x),\quad \mathrm{pair}\,a\,b=\lambda\gamma.(a\gamma,b\gamma),\\
&\mathrm{fst}\,s=\lambda\gamma.\pi_1(s\gamma),\quad \mathrm{snd}\,s=\lambda\gamma.\pi_2(s\gamma),\\
&\text{all equations}=\mathrm{refl}
\;]
\end{array}
\]

\section{Definition of a parametricity model}
\label{twoltt-parametricity-model-def}

\[\begin{array}{ll}
    \mathrm{Para}_\ell = &[\;
\mathcal C=(\Gamma_0\ \Gamma_1:\mathcal U_{\ell+1})\times(\Gamma_0\to\Gamma_1\to\mathcal U_{\ell+1}),\\
&(\Gamma_0,\Gamma_1,\Gamma_r)\triangleright(\Delta_0,\Delta_1,\Delta_r)=\\ &\qquad=(f_0:\Gamma_0\to\Delta_0)\times(f_1:\Gamma_1\to\Delta_1)\times
 \big((\gamma_0\,\gamma_1)\to\Gamma_r\,\gamma_0\,\gamma_1\to\Delta_r\,(f_0\gamma_0)\,(f_1\gamma_1)\big),\\
&\mathrm{id}=(\lambda x.x,\lambda x.x,\lambda\_\_\,r.r),\\
&(f_0,f_1,f_r)\circ(g_0,g_1,g_r)=(f_0\circ g_0,\ f_1\circ g_1,\ \lambda\gamma_0\gamma_1 r.f_r\,\_\,\_\,(g_r\,\gamma_0\,\gamma_1\,r)),\\
&\cdot=(\top,\top,\lambda\_\_.\top),\\
&\mathrm{Ty}_j\,(\Gamma_0,\Gamma_1,\Gamma_r)=\\&\qquad=
(A_0:\Gamma_0\to\mathcal U_j)\times(A_1:\Gamma_1\to\mathcal U_j)\times\\&\qquad\quad\times
 \big((\gamma_0\,\gamma_1)(\gamma_r:\Gamma_r\gamma_0\gamma_1)\to A_0\gamma_0\to A_1\gamma_1\to\mathcal U_j\big),\\
&\mathrm{Tm}_j\,\Gamma\,(A_0,A_1,A_r)=\\&\qquad=(t_0:(\gamma_0:\Gamma_0)\to A_0\gamma_0)\times(t_1:(\gamma_1:\Gamma_1)\to A_1\gamma_1)\times\\&\qquad\quad\times
 \big((\gamma_0\,\gamma_1\,\gamma_r)\to A_r\,\gamma_0\gamma_1\gamma_r\,(t_0\gamma_0)\,(t_1\gamma_1)\big),\\
&\Gamma.A=\big((\gamma_0:\Gamma_0)\times A_0\gamma_0,\ (\gamma_1:\Gamma_1)\times A_1\gamma_1,\\&\qquad
 \lambda(\gamma_0,a_0)(\gamma_1,a_1).(\gamma_r:\Gamma_r\gamma_0\gamma_1)\times A_r\,\gamma_0\gamma_1\gamma_r\,a_0\,a_1\big),\\
&p=(\pi_1,\pi_1,\lambda\_\_.\pi_1),\quad q=(\pi_2,\pi_2,\lambda\_\_.\pi_2),\\
&\langle\sigma,t\rangle=(\lambda\delta.(\sigma_0\delta,t_0\delta),\ \lambda\delta.(\sigma_1\delta,t_1\delta),\
 \lambda\delta_0\delta_1\delta_r.(\sigma_r\,\delta_0\delta_1\delta_r,\ t_r\,\delta_0\delta_1\delta_r)),\\
&\mathrm U_j=(\lambda\_.\mathcal U_j,\ \lambda\_.\mathcal U_j,\ \lambda\_\_\_\,A_0\,A_1.\,A_0\to A_1\to\mathcal U_j),\\
&\mathrm{El}\,(t_0,t_1,t_r)=(\lambda\gamma_0.\mathrm{El}(t_0\gamma_0),\ \lambda\gamma_1.\mathrm{El}(t_1\gamma_1),\ t_r),\\
&\mathrm{Nat}_j=(\lambda\_.\mathrm{Nat}_j,\ \lambda\_.\mathrm{Nat}_j,\ \lambda\_\_\_\,m\,n.\,m=_{\mathrm{Nat}_j}n),\\
&0=(\lambda\_.0,\lambda\_.0,\lambda\_\_\_.\mathrm{refl}),\quad
\mathrm S\,t=(\lambda\gamma.\mathrm S(t_0\gamma),\ \lambda\gamma.\mathrm S(t_1\gamma),\ \lambda\gamma_0\gamma_1\gamma_r.\mathrm{ap}_{\mathrm S}(t_r\,\gamma_0\gamma_1\gamma_r)),\\
&\Pi\,A\,B=\big(\lambda\gamma_0.(x_0:A_0\gamma_0)\to B_0(\gamma_0,x_0),\ \lambda\gamma_1.(x_1:A_1\gamma_1)\to B_1(\gamma_1,x_1),\\
&\qquad\quad \lambda\gamma_0\gamma_1\gamma_r\,f_0\,f_1.(x_0\,x_1)(x_r:A_r\,\gamma_0\gamma_1\gamma_r\,x_0\,x_1)\to\\&\qquad\to B_r\,(\gamma_0,x_0)(\gamma_1,x_1)(\gamma_r,x_r)\,(f_0x_0)\,(f_1x_1)\big),\\
&\mathrm{lam}\,b=(\lambda\gamma_0x_0.b_0(\gamma_0,x_0),\ \lambda\gamma_1x_1.b_1(\gamma_1,x_1),\ \lambda\gamma_0\gamma_1\gamma_rx_0x_1x_r.b_r(\gamma_0,x_0)(\gamma_1,x_1)(\gamma_r,x_r)),\\
&\mathrm{app}\,f\,a=(\lambda\gamma_0.f_0\gamma_0(a_0\gamma_0),\ \lambda\gamma_1.f_1\gamma_1(a_1\gamma_1),\ \lambda\gamma_0\gamma_1\gamma_r.f_r\,\gamma_0\gamma_1\gamma_r\,(a_0\gamma_0)(a_1\gamma_1)(a_r\gamma_0\gamma_1\gamma_r)),\\
&\Sigma\,A\,B=\big(\lambda\gamma_0.(x_0:A_0\gamma_0)\times B_0(\gamma_0,x_0),\ \dots,\
 \lambda\gamma_0\gamma_1\gamma_r\,s_0\,s_1.\\
 &\qquad\quad(x_r:A_r\,\gamma_0\gamma_1\gamma_r\,(\pi_1s_0)(\pi_1s_1))\times B_r\,\_\,\_\,(\gamma_r,x_r)\,(\pi_2s_0)(\pi_2s_1)\big),\\
&\mathrm{Id}\,A\,t\,t'=\big(\lambda\gamma_0.t_0\gamma_0=t_0'\gamma_0,\ \lambda\gamma_1.t_1\gamma_1=t_1'\gamma_1,\
 \\&\qquad\lambda\gamma_0\gamma_1\gamma_r\,e_0\,e_1.\ \mathrm{tr}^{A_r\gamma_0\gamma_1\gamma_r}_{e_0,e_1}(t_r\,\gamma_0\gamma_1\gamma_r)=t'_r\,\gamma_0\gamma_1\gamma_r\big),\\
&\mathrm{refl}\,e=(\dots,\mathrm{refl},\mathrm{refl}),\\
&\mathrm u_j,\ \mathrm{nat}_j,\ \mathrm{pi},\ \mathrm{sigma},\ \mathrm{id}_j,\ \mathrm{natind},\ \mathrm J:\ \text{componentwise},\\
&\text{all equations}=\mathrm{refl}
\;]
\end{array}
\]


\end{document}
